% ============================================================================
% RELATÓRIO FINAL DE ATIVIDADES — PIBIC/UFPE, edição 2025-2026
%
% Segue o modelo "Relatório final de atividades do estudante bolsista e
% voluntário dos programas PIBIC, PIBIC-AF, PIBITI e PIBIC-EM/CNPq - UFPE"
% publicado pela Propesqi em www.ufpe.br/iniciacao-cientifica (modelo mais
% recente disponível, 2022-23): Times New Roman 12, espaçamento simples,
% páginas numeradas, seções 1 a 11 na ordem do modelo, resumo com no máximo
% 200 palavras, introdução com no máximo 3 páginas, resultados com no máximo
% 8 páginas, documento com no máximo 25 páginas e 3 MB, enviado em PDF.
%
% Compilar:  docs/relatorio_ic/build.sh   (-> build/relatorio_final_pibic.pdf)
%
% Campos de identificação já confirmados. O comando \pendente{...} (destaque em
% amarelo) fica disponível para marcar qualquer trecho que ainda precise de revisão.
% ============================================================================
\documentclass[12pt,a4paper]{article}

% --- codificação, idioma e fonte ------------------------------------------
\usepackage[utf8]{inputenc}
\usepackage[T1]{fontenc}
\usepackage[brazil]{babel}
\usepackage{amsmath}
\usepackage{newtxtext,newtxmath}      % Times New Roman (texto e matemática)

% --- página ----------------------------------------------------------------
\usepackage[a4paper,margin=2.5cm]{geometry}   % espaçamento simples = padrão do article
\setlength{\parindent}{1.25cm}
\setlength{\parskip}{2pt plus 1pt}
\usepackage{fancyhdr}
\pagestyle{fancy}
\fancyhf{}
\fancyhead[R]{\thepage}
\renewcommand{\headrulewidth}{0pt}
\setlength{\headheight}{15pt}

% --- títulos das seções (numeradas como no modelo: "1. IDENTIFICAÇÃO") -----
\usepackage{titlesec}
\titleformat{\section}{\normalfont\bfseries}{\thesection.}{0.5em}{}
\titleformat{\subsection}{\normalfont\bfseries}{\thesubsection}{0.5em}{}
\titleformat{\subsubsection}{\normalfont\bfseries\itshape}{\thesubsubsection}{0.5em}{}
\titlespacing*{\section}{0pt}{11pt plus 3pt minus 2pt}{5pt}
\titlespacing*{\subsection}{0pt}{8pt plus 2pt minus 1pt}{3pt}
\titlespacing*{\subsubsection}{0pt}{6pt plus 2pt minus 1pt}{2pt}
\setcounter{tocdepth}{1}
\usepackage[titles]{tocloft}            % sumário compacto
\setlength{\cftbeforesecskip}{2pt}
\setlength{\cftaftertoctitleskip}{6pt}

% --- figuras, tabelas e listas --------------------------------------------
\usepackage{graphicx,url}
\usepackage{booktabs,multirow,array}
\usepackage{xcolor}
\usepackage{tikz}
\usetikzlibrary{shapes.geometric, arrows.meta, positioning, calc, fit, backgrounds}
\usepackage{pgfplots}
\pgfplotsset{compat=1.18}
\usepackage{caption}
\captionsetup{font=small,labelfont=bf,labelsep=period}
\usepackage{enumitem}
\setlist[itemize]{noitemsep,topsep=2pt,leftmargin=*}
\setlist[enumerate]{noitemsep,topsep=2pt,leftmargin=*}

% bibliografia como seção numerada, com o título do modelo
\usepackage[nottoc,numbib]{tocbibind}
\addto\captionsbrazil{\renewcommand{\refname}{REFERÊNCIAS BIBLIOGRÁFICAS}}
% menos espaço entre as entradas da bibliografia
\let\oldthebibliography\thebibliography
\renewcommand{\thebibliography}[1]{\oldthebibliography{#1}\setlength{\itemsep}{2pt plus 1pt}}

% permite quebra de linha em URLs longas (notas de rodapé com links do GitHub)
\expandafter\def\expandafter\UrlBreaks\expandafter{\UrlBreaks%
  \do\a\do\b\do\c\do\d\do\e\do\f\do\g\do\h\do\i\do\j\do\k\do\l\do\m%
  \do\n\do\o\do\p\do\q\do\r\do\s\do\t\do\u\do\v\do\w\do\x\do\y\do\z%
  \do\A\do\B\do\C\do\D\do\E\do\F\do\G\do\H\do\I\do\J\do\K\do\L\do\M%
  \do\N\do\O\do\P\do\Q\do\R\do\S\do\T\do\U\do\V\do\W\do\X\do\Y\do\Z%
  \do\-\do\_}

\usepackage[hidelinks]{hyperref}

% campo ainda não preenchido: destacado em amarelo no PDF
\newcommand{\pendente}[1]{\colorbox{yellow!55}{\strut #1}}

\sloppy

\begin{document}

% ============================================================================
% CABEÇALHO (idêntico ao modelo da Propesqi)
% ============================================================================
\begin{center}
\includegraphics[width=5.4cm]{logo_ufpe_propesqi.jpeg}\\[12pt]
{\bfseries RELATÓRIO FINAL DE ATIVIDADES DO ESTUDANTE\\
BOLSISTA E VOLUNTÁRIO DOS PROGRAMAS\\
PIBIC, PIBIC-AF, PIBITI E PIBIC-EM/CNPq -- UFPE}\\[2pt]
(Refere-se às atividades realizadas no período de setembro de 2025 a agosto de 2026)
\end{center}

% ============================================================================
\section{IDENTIFICAÇÃO}
\label{sec:identificacao}

\noindent
\begin{tabular}{|p{5.9cm}|p{9.1cm}|}
\hline
Nome do(a) Orientador(a): &
  Andson M. Balieiro (orientador) e Caio B. B. de Souza (coorientador) \\ \hline
Nome do(a) Estudante: &
  Luiz Henrique B. A. da Silva \\ \hline
Área do projeto: &
  Ciência da Computação (Redes de Computadores e Segurança da Informação) \\ \hline
Título do projeto: &
  Tecnologias e Mecanismos Sustentáveis para Segurança e Alocação de Recursos
  em Redes 5G Avançadas e 6G (plano de trabalho: Uma Solução Hierárquica para
  Detecção de Ciberataques em Ambientes IoT baseados em Redes MEC 6G) \\ \hline
Houve alteração do nome do projeto original? Se sim, acrescente aqui o nome do
projeto originalmente aprovado &
  Não \\ \hline
ID do projeto &
  PIC2570-2025 \\ \hline
ODS na qual projeto se insere &
  ODS 9 -- Indústria, Inovação e Infraestrutura: construir infraestruturas
  resilientes, promover a industrialização inclusiva e sustentável e fomentar
  a inovação \\ \hline
\end{tabular}

% ============================================================================
\section{RESUMO}
\label{sec:resumo}

A Internet das Coisas (IoT) e as redes de sexta geração (6G) ampliam a
superfície de ataque em ambientes de Computação de Borda de Acesso Múltiplo
(MEC), cujos dispositivos têm recursos limitados. Este projeto implementou e
validou um Sistema de Detecção de Intrusão (IDS) hierárquico de
duas fases para ambientes IoT/MEC, executado em um nó de borda ARM (VIM~4). A
Fase~1 é um classificador binário LightGBM otimizado pelo escore F2; a Fase~2,
acionada só sobre fluxos suspeitos, classifica o tipo de ataque com LightGBM
multiclasse. O pipeline de dados converte capturas PCAP de dois datasets do
CIC, com a ferramenta \textit{netflower}, criada no projeto, em um banco SQLite
de $\sim$190~GB. Em avaliação \textit{offline}, a Fase~1 alcançou revocação
de 96\% na classe de ataque (F2 = 88,6\%). No \textit{testbed}, com oito
categorias de ataque, a detecção é avaliada \emph{por fluxo} (cada fluxo
rotulado pela janela de ataque em que começou, métrica imune ao atraso de
classificação): o IDS binário alcançou precisão por fluxo de $\approx$98\% e
revocação de $\approx$64\% (F1 de $\approx$78\%), mais forte nas inundações
volumétricas de sessão TCP (\textit{spoofing} com revocação de 94\% e DoS de
92\%); a revocação reflete o limiar de 0,9, que descarta os \textit{floods} UDP de
pacote único. A Fase~2 classificou o reconhecimento de rede com F1 = 73,7\%, mas
colapsa os demais tipos na classe DoS (macro-F1 de 19,4\% sobre as classes de
ataque). A validação mede ainda um teto de vazão do nó de borda: a $\sim$1.050
fluxos/s de \textit{flood} o VIM~4 não classifica em tempo real, enfileirando o
tráfego com latência mediana de 180 a 437~s, e a segunda fase leva a CPU de 12\% a
40\% e a potência de 3,5 a 5,1~W. Os resultados caracterizam onde as
\textit{features} por fluxo deixam de separar tipos de ataque e o limite de
capacidade do dispositivo sob ataque volumétrico.

\renewcommand{\contentsname}{SUMÁRIO}
\tableofcontents

% ============================================================================
\section{INTRODUÇÃO}
\label{sec:introducao}

A convergência entre redes móveis da Sexta Geração (6G), Internet das Coisas
(IoT) e Computação de Borda de Acesso Múltiplo (MEC) promete transformar a
infraestrutura de comunicação digital, com estimativas de densidades de
dispositivos IoT na ordem de $10^6$ por km² e latências de comunicação
inferiores a 1~ms, sendo a tecnologia MEC um componente essencial para
viabilizar esses requisitos~\cite{you:21}. Contudo, a massificação de
dispositivos expande a superfície de ataques cibernéticos: dispositivos IoT
frequentemente possuem poder computacional e energia reduzidos, o que
inviabiliza a adoção de soluções de segurança tradicionais diretamente nos
terminais.

Na literatura, Sistemas de Detecção de Intrusão (IDS) baseados em regras e
assinaturas têm sido implantados em nós de borda (MEC) para aproximar a
detecção dos dispositivos IoT e reduzir a latência de resposta. Essas soluções
comparam o tráfego observado a padrões de ataque previamente catalogados e,
por isso, alcançam alta precisão sobre ameaças conhecidas~\cite{lazarevic:05}.
Sua limitação estrutural está no que não conhecem: ataques inéditos ou
variantes de ataques existentes não correspondem a nenhuma assinatura
registrada e atravessam o detector sem gerar alerta. Além disso, manter o
conjunto de regras atualizado exige intervenção manual constante, o que se
torna impraticável na escala e na diversidade de dispositivos previstas para
as redes 6G. Diante dessas limitações, IDS baseados em aprendizado de máquina
executados em nós de borda surgem como alternativa viável: treinados com dados
históricos de tráfego, classificam fluxos de rede em tempo real e generalizam
para padrões de ataque não vistos, mantendo baixo custo computacional na
inferência~\cite{gyamfi:22}. A combinação da detecção por assinatura com a
detecção por anomalia em arquiteturas hierárquicas representa o estado da arte
para ambientes com restrições de recursos, como dispositivos IoT e nós de
borda MEC~\cite{zolanvari:19}.

Trabalhos recentes corroboram essa direção. Os datasets do Canadian Institute
for Cybersecurity (CIC) são amplamente usados para treinar e avaliar IDS em
redes IoT; o CIC IoT Dataset 2023~\cite{neto:23}, por exemplo, reproduz uma
topologia de 105 dispositivos com 33 ataques organizados em sete categorias:
Negação de Serviço Distribuída (DDoS), Negação de Serviço (DoS),
\textit{recon}, ataques web, força bruta, \textit{spoofing} e Mirai. Sobre
essa base, Houichi et~al.~\cite{houichi:25} comparam classificadores clássicos
de aprendizado de máquina para segurança de cidades inteligentes, e Chen
et~al.~\cite{chen:24} combinam LightGBM com seleção de \textit{features} por um
algoritmo \textit{salp swarm} aprimorado; ambos, porém, reportam apenas
métricas \textit{offline} sobre o conjunto de teste do próprio dataset, sem
avaliar o modelo implantado em um nó de borda real nem o desvio entre o
tráfego do dataset e o de produção. Mover a detecção para a borda é motivado
pela baixa latência e pelo volume de tráfego em IoT: Gyamfi e
Jurcut~\cite{gyamfi:22} revisam abordagens de projeto de IDS que combinam MEC,
aprendizado de máquina e datasets, e Yang et~al.~\cite{yang:23} propõem
detecção eficiente por colaboração \textit{cloud--edge}. Nessas arquiteturas,
contudo, a etapa com maior carga de processamento é tipicamente movida para a
nuvem central, e o custo de manter a detecção no próprio nó de borda, em
termos de CPU, RAM e energia, raramente é medido. A ideia de filtrar o tráfego
com um classificador leve antes de acionar um mais custoso também é recorrente:
Aldaej et~al.~\cite{aldaej:24} propõem um IDS de dois estágios para
plataformas IoT-borda, com um classificador binário inicial seguido de um
\textit{ensemble} sobre os fluxos sinalizados. A presente proposta compartilha
essa filosofia, mas distingue-se por executar ambas as fases no mesmo nó de
borda e por isolar e medir o custo marginal da segunda fase diretamente no
dispositivo. No horizonte das redes 6G, densidade de dispositivos e metas de
latência sub-milissegundo sem precedentes~\cite{you:21} expandem a superfície
de ataque ao mesmo tempo em que reduzem a janela disponível para detectar e
reagir a uma intrusão; nessa direção, Chinnasamy et~al.~\cite{chinnasamy:25}
discutem sistemas de detecção e prevenção de intrusão guiados por inteligência
artificial para redes 6G e reforçam a necessidade de detectores que conciliem
alta acurácia e baixo custo computacional em nós de borda, exatamente o regime
de operação avaliado neste trabalho. A Tabela~\ref{tab:related} sintetiza esse
posicionamento segundo as cinco dimensões que distinguem a presente proposta.

\begin{table}[ht]
\centering
\caption{Posicionamento do trabalho frente à literatura segundo cinco dimensões de avaliação.}
\label{tab:related}
\setlength{\tabcolsep}{4pt}
\resizebox{\textwidth}{!}{%
\begin{tabular}{@{}lccccc@{}}
\toprule
\textbf{Trabalho}
  & \shortstack{\textbf{Deploy em}\\\textbf{borda real}}
  & \shortstack{\textbf{Avaliação}\\\textbf{online}}
  & \shortstack{\textbf{Custo no}\\\textbf{dispositivo}}
  & \shortstack{\textbf{Multi-}\\\textbf{estágio}}
  & \shortstack{\textbf{Foco}\\\textbf{6G}}
\\ \midrule
Houichi et~al.~\cite{houichi:25}       & Não & Não & Não & Não & Não \\
Chen et~al.~\cite{chen:24}             & Não & Não & Não & Não & Não \\
Gyamfi e Jurcut~\cite{gyamfi:22}       & \multicolumn{5}{c}{\textit{revisão de literatura}} \\
Yang et~al.~\cite{yang:23}             & Parcial & Não & Não & Sim & Não \\
Aldaej et~al.~\cite{aldaej:24}         & Sim & Não & Não & Sim & Não \\
Chinnasamy et~al.~\cite{chinnasamy:25} & Não & Não & Não & Não & Sim \\
\textbf{Este trabalho}                 & \textbf{Sim} & \textbf{Sim} & \textbf{Sim} & \textbf{Sim} & \textbf{Sim} \\
\bottomrule
\end{tabular}}
\end{table}

Esse conjunto de trabalhos avança em modelos de detecção, em arquiteturas de
borda e nos requisitos de segurança de 6G, mas nenhum deles acompanha um IDS
hierárquico do treinamento à operação em um dispositivo de borda real. Somam-se
a essa lacuna dois desafios em aberto: conjuntos de dados capturados em
ambientes controlados induzem desvio de distribuição (\textit{distribution
shift}) quando o modelo encontra tráfego real, degradando a acurácia; e a
diversidade de tipos de ataque demanda objetivos de otimização distintos, como
latência mínima na borda e maior capacidade de detecção nas camadas
superiores, o que motiva uma abordagem hierárquica em que cada camada da
arquitetura MEC é otimizada para um objetivo próprio.

Este projeto endereça essas lacunas propondo, implementando e validando um IDS
hierárquico de duas fases para ambientes IoT/MEC. A primeira fase é um
classificador binário leve, executado no nó de borda, que rotula cada fluxo de
rede como benigno ou malicioso. A segunda é um classificador multiclasse,
acionado apenas sobre os fluxos sinalizados como maliciosos, que os atribui a
classes específicas de ataque. O trabalho cobre o pipeline completo, desde a
coleta de dados em formato PCAP até a implantação e a avaliação da solução em
um dispositivo embarcado (Khadas VIM~4), com tráfego de ataque real. Cabe
delimitar o sentido em que os termos MEC e 6G se aplicam: a avaliação não
ocorre sobre uma rede 6G operacional, cuja padronização ainda está em
curso~\cite{you:21}, nem sobre uma infraestrutura MEC comercial completa. O
\textit{testbed} instancia o elemento dessa arquitetura no qual o IDS reside,
um nó de computação de borda ARM posicionado junto aos dispositivos que
protege, no sentido do paradigma de computação na borda de acesso
múltiplo~\cite{gyamfi:22}, cujo requisito definidor, inspecionar o tráfego
próximo aos dispositivos dentro de um orçamento restrito de recursos, é
exatamente o que o nó avaliado reproduz. Já 6G designa o regime de requisitos
que motiva e dimensiona a solução. Como o IDS opera sobre fluxos IP
reconstruídos acima da camada de acesso, a detecção independe da tecnologia de
rádio subjacente; os resultados obtidos na rede local do \textit{testbed}
caracterizam o custo e o comportamento do detector no nó de borda e
transferem-se a implantações MEC em redes 6G, nas quais muda o enlace de
acesso, não a representação de fluxos sobre a qual o detector decide.

As principais contribuições são um pipeline de dados reprodutível baseado em
SQLite para amostragem balanceada em escala; a ferramenta \textit{netflower},
publicada no PyPI, para conversão paralela de PCAPs em fluxos e captura ao
vivo em dispositivos de borda; o IDS hierárquico com duas fases
complementares; e a validação no \textit{testbed} com tráfego de ataque real,
com comparação quantitativa entre o pipeline binário e o multiclasse em
métricas de detecção (acurácia, precisão, revocação, F1 e taxas de falsos
positivos, FPR, e de falsos negativos, FNR) e de custo (CPU, RAM, energia e
\textit{throughput} de alertas). Nos resultados mais relevantes, avaliados
\emph{por fluxo} (cada fluxo rotulado pela janela em que começou, métrica imune
ao atraso de classificação), o IDS binário alcançou precisão de $\approx$98\% e
revocação de $\approx$64\% (F1 de $\approx$78\%), forte nas inundações
volumétricas de sessão TCP (\textit{spoofing} com revocação de 94\% e DoS de
92\%), com a revocação limitada pelo limiar de 0,9, que descarta os \textit{floods}
UDP de pacote único. A Fase~2 classificou corretamente o reconhecimento de rede
(F1 = 73,7\%) e delimitou empiricamente onde as 55 \textit{features} por fluxo
deixam de separar os tipos de ataque, colapsando os demais (\textit{spoofing},
força bruta, \textit{Man-in-the-Middle}, MITM) na classe DoS (macro-F1 de 19,4\%).
A validação mede ainda o teto de vazão do dispositivo: sob \textit{flood} o VIM~4
não classifica em tempo real e enfileira o tráfego com minutos de latência, e a
segunda fase custa CPU e energia mensuráveis.
As seções seguintes apresentam os objetivos do plano de trabalho
(Seção~\ref{sec:objetivos}), a metodologia (Seção~\ref{sec:metodologia}), os
resultados da validação no \textit{testbed} (Seção~\ref{sec:resultados}) e as
conclusões (Seção~\ref{sec:conclusoes}).

% ============================================================================
\section{OBJETIVOS}
\label{sec:objetivos}

O objetivo geral do plano de trabalho foi propor, implementar e validar
experimentalmente um IDS hierárquico para ambientes IoT/MEC em redes 6G,
capaz de executar integralmente em um nó de borda ARM de recursos limitados,
conciliando alta capacidade de detecção com baixo custo de CPU, memória e
energia.

Os objetivos específicos foram:

\begin{enumerate}
  \item construir um pipeline de dados reprodutível que converta capturas
    brutas (PCAP) dos datasets CIC IoT 2023 e CIC IIoT 2025, complementadas
    por tráfego benigno real, em uma base de fluxos consultável e amostrável
    em escala (SQLite);
  \item desenvolver uma ferramenta de extração de fluxos bidirecionais
    compatível com o conjunto de \textit{features} do CICFlowMeter e leve o
    suficiente para captura ao vivo em dispositivos ARM (\textit{netflower});
  \item definir a taxonomia de ataques e o conjunto final de
    \textit{features}, removendo atributos que induzem memorização de hosts
    ou desvio entre treino e implantação;
  \item treinar e otimizar a Fase~1 (classificador binário, escore F2) e a
    Fase~2 (classificador multiclasse, macro F1) com LightGBM e Optuna, e
    avaliar seu desempenho \textit{offline};
  \item implantar as duas fases no nó de borda (Khadas VIM~4) e validá-las em
    \textit{testbed} com tráfego real das oito categorias de ataque, medindo
    métricas de detecção (precisão, revocação, F1, FPR e FNR) e de custo (CPU,
    RAM, energia e \textit{throughput});
  \item comparar quantitativamente o pipeline binário (Fase~1) e o
    hierárquico completo (Fases~1+2), isolando o custo marginal da segunda
    fase no dispositivo;
  \item explorar, de forma preliminar e \textit{offline}, uma terceira fase de
    agrupamento não supervisionado para os fluxos classificados com baixa
    confiança pela Fase~2, como base para a detecção de ameaças desconhecidas
    em trabalhos futuros.
\end{enumerate}

% ============================================================================
\section{METODOLOGIA}
\label{sec:metodologia}

A metodologia adotada cobre todo o ciclo do sistema, desde a coleta de capturas
brutas de rede até a avaliação do detector em operação em um nó de borda real.
A Figura~\ref{fig:metodologia} apresenta essa visão geral e o encadeamento entre
as etapas. O tráfego é coletado em arquivos de captura de pacotes (\textit{Packet
Capture}, PCAP), convertido em fluxos bidirecionais pela ferramenta
\textit{netflower} e consolidado em um banco SQLite. Sobre esse banco aplica-se a
engenharia de características que reduz o vetor de \textit{features}. Os dados
resultantes treinam os dois classificadores que compõem o IDS hierárquico e os
modelos treinados são implantados no nó de borda (placa VIM~4) para validação no
\textit{testbed} com tráfego de ataque real. As subseções seguintes detalham cada
etapa na ordem da Figura~\ref{fig:metodologia}: a arquitetura hierárquica
(Seção~\ref{subsec:arquitetura}), o pipeline de dados
(Seção~\ref{subsec:pipeline}), o mapeamento de rótulos
(Seção~\ref{subsec:rotulos}), a engenharia de características
(Seção~\ref{subsec:features}), o treinamento das duas fases
(Seção~\ref{subsec:treinamento}), o protocolo de validação no \textit{testbed}
(Seção~\ref{subsec:validacao}) e as métricas usadas na avaliação
(Seção~\ref{subsec:metricas}).

\begin{figure}[ht]
\centering
\resizebox{\textwidth}{!}{%
\begin{tikzpicture}[
  font=\footnotesize,
  stage/.style={rectangle, rounded corners, draw, thick, align=center,
                minimum height=1.15cm, minimum width=2.0cm, fill=gray!8},
  arr/.style={-{Stealth[length=2.5mm]}, thick}]
  \node[stage] (coleta)  {Coleta\\(PCAPs CIC +\\benigno real)};
  \node[stage, right=0.7cm of coleta] (flow)  {Extração de\\fluxos\\(\textit{netflower})};
  \node[stage, right=0.7cm of flow]   (db)    {Banco\\SQLite\\($\sim$190~GB)};
  \node[stage, right=0.7cm of db]     (feat)  {Engenharia de\\características\\(83\,$\rightarrow$\,55)};
  \node[stage, right=0.7cm of feat]   (train) {Treinamento\\(Optuna)\\Fase~1 + Fase~2};
  \node[stage, right=0.7cm of train]  (deploy){Implantação\\no nó de borda\\(VIM~4)};
  \node[stage, right=0.7cm of deploy] (val)   {Validação no\\\textit{testbed} (ataques\\+ métricas)};
  \draw[arr] (coleta) -- (flow);
  \draw[arr] (flow)   -- (db);
  \draw[arr] (db)     -- (feat);
  \draw[arr] (feat)   -- (train);
  \draw[arr] (train)  -- (deploy);
  \draw[arr] (deploy) -- (val);
\end{tikzpicture}}
\caption{Visão geral da metodologia: da coleta de PCAPs à validação no \textit{testbed} no nó de borda.}
\label{fig:metodologia}
\end{figure}

\subsection{Arquitetura hierárquica de duas fases}
\label{subsec:arquitetura}

Devido às restrições inerentes a ambientes IoT/MEC, adotou-se uma arquitetura
hierárquica de duas fases. Como a maioria do tráfego em qualquer rede é benigna,
aplicar um classificador pesado sobre todos os fluxos desperdiçaria energia e
capacidade do nó de borda. Além disso, diferentes etapas do pipeline têm
objetivos de otimização distintos: a detecção binária privilegia a revocação,
para não deixar ataques passarem, enquanto a classificação de tipo privilegia a
precisão por classe. Um pipeline progressivo concilia essas restrições, pois a
Fase~1 filtra o tráfego benigno com custo mínimo e somente os fluxos suspeitos
avançam à Fase~2 para classificação de tipo.

A arquitetura, ilustrada na Figura~\ref{fig:arquitetura}, distribui a carga
computacional hierarquicamente em duas fases complementares:

\begin{enumerate}
  \item \textbf{Fase~1, classificador binário (nó de borda):} decide
    \textit{benigno vs.\ ataque} com mínimo custo computacional, atuando como um
    filtro de primeira linha sobre todo o tráfego observado na interface de rede.
  \item \textbf{Fase~2, classificador multiclasse (nó de borda ou host):}
    identifica o tipo específico de ataque apenas para os fluxos sinalizados pela
    Fase~1. Por ser mais pesada, é a candidata natural a descarregamento para um
    host dedicado quando o nó de borda for muito restrito.
\end{enumerate}

O encadeamento das duas fases tem efeito direto sobre o custo de operação. Cada
fluxo extraído pelo \textit{netflower} é primeiro submetido à Fase~1. Caso a
probabilidade de ataque estimada não atinja o limiar de decisão, o fluxo é
tratado como benigno e descartado sem custo adicional, de modo que o volume
dominante de tráfego legítimo nunca chega ao segundo modelo. Apenas a fração
sinalizada como suspeita aciona a Fase~2, que lhe atribui um rótulo de tipo.
Como a inferência da Fase~2 ocorre em série após a Fase~1 e sobre um número
reduzido de fluxos, seu custo marginal é pequeno e amortizado pelo rastreamento
de fluxos, que já domina o tempo de processamento
(Seção~\ref{sec:resultados}). Neste trabalho, ambas as fases foram avaliadas
executando no próprio nó de borda, um dispositivo VIM~4, enquanto um computador
pessoal (PC) atuou como gerador de ataques e como plataforma de treinamento dos
modelos. Manter a Fase~2 no nó de borda, em vez de descarregá-la em um host
separado, permite medir o custo adicional do segundo modelo diretamente no
dispositivo de borda (Seção~\ref{sec:resultados}).

\begin{figure}[ht]
\centering
\begin{tikzpicture}[
  font=\footnotesize,
  box/.style={rectangle, rounded corners, draw, thick, align=center,
              minimum height=1.0cm, minimum width=3.0cm, fill=gray!8},
  dec/.style={diamond, aspect=2, draw, thick, align=center, fill=gray!12,
              inner sep=1pt},
  term/.style={rectangle, draw, thick, align=center, minimum height=0.9cm,
               minimum width=2.4cm},
  arr/.style={-{Stealth[length=2.5mm]}, thick}]
  \node[term] (traf) {Tráfego de rede\\(\texttt{eth0})};
  \node[box, below=0.8cm of traf] (f1) {\textbf{Fase~1}\\LightGBM binário};
  \node[dec, below=0.8cm of f1] (dec) {$P_1 \geq 0{,}9$?};
  \node[term, left=1.6cm of dec] (ben) {Benigno\\(descartado)};
  \node[box, below=0.9cm of dec] (f2) {\textbf{Fase~2}\\LightGBM multiclasse};
  \node[term, below=0.8cm of f2] (tipo) {Alerta tipado\\(\texttt{dos}, \texttt{recon}, \ldots)};
  \draw[arr] (traf) -- (f1);
  \draw[arr] (f1) -- (dec);
  \draw[arr] (dec) -- node[above]{não} (ben);
  \draw[arr] (dec) -- node[right]{sim} (f2);
  \draw[arr] (f2) -- (tipo);
  \begin{scope}[on background layer]
    \node[draw, dashed, rounded corners, fit=(traf)(f1)(dec)(f2)(tipo)(ben),
          inner sep=0.35cm, label={[anchor=south, font=\small, yshift=-2pt]north:nó de borda (VIM~4)}] {};
  \end{scope}
\end{tikzpicture}
\caption{Arquitetura hierárquica de duas fases executando no nó de borda (VIM~4).}
\label{fig:arquitetura}
\end{figure}

\subsection{Pipeline de dados}
\label{subsec:pipeline}

Para treinar os modelos, foi definido um pipeline de dados que converte capturas
brutas de rede em uma base consultável e amostrável em escala e que, portanto,
condiciona a qualidade de toda a detecção. O pipeline articula três elementos
que operam em sequência: os \textit{datasets} de origem, que combinam tráfego de
ataque público com tráfego benigno real coletado pelo estudante; a ferramenta
\textit{netflower}, responsável por converter as capturas em fluxos
bidirecionais com \textit{features} estatísticas; e o banco SQLite, que
consolida esses fluxos para consulta e amostragem.

Para compor a base, apenas datasets com PCAPs separados por categoria foram
utilizados, nos quais os rótulos derivam do nome das pastas (\texttt{benign/},
\texttt{ddos/}, \texttt{malware/} etc.), evitando inconsistências de esquema
entre fontes~\cite{sharafaldin:18}. Foram empregados dois datasets do
CIC: o CIC IoT Dataset
2023~\cite{neto:23}\footnote{\url{https://www.unb.ca/cic/datasets/iotdataset-2023.html}} e o
CIC IIoT Dataset 2025\footnote{\url{https://www.unb.ca/cic/datasets/iiot-dataset-2025.html}},
descritos na Tabela~\ref{tab:datasets}.

\begin{table}[ht]
\centering
\caption{Datasets CIC utilizados no pipeline de dados.}
\label{tab:datasets}
\begin{tabular}{@{}ll@{}}
\toprule
\textbf{Dataset} & \textbf{Foco principal} \\ \midrule
CIC IoT Dataset 2023 & Tráfego multi-ataque em dispositivos IoT \\
CIC IIoT 2025        & Intrusões em IoT industrial \\
\bottomrule
\end{tabular}
\end{table}

Além do tráfego benigno proveniente dos datasets públicos, a classe benigna foi
complementada com tráfego benigno real capturado no ambiente experimental. Ao
longo de dezenas de horas, um terminal registrou continuamente os PCAPs do uso
cotidiano de um computador pessoal, incluindo tráfego de aplicações de estudo,
navegação web, \textit{streaming} de vídeo, jogos, sessões SSH e atualizações
de sistema, posteriormente convertidos em fluxos pelo \textit{netflower} e
rotulados como \texttt{benign}. Esse material atua como um terceiro conjunto de
tráfego benigno, com a vantagem de refletir o comportamento real de um host em
uso interativo e não tráfego sintético. A motivação é metodológica: o tráfego
benigno dos datasets CIC é gerado de forma \textit{scripted}, com padrões
temporais regulares, ao passo que o tráfego interativo real apresenta
características distintas, como \textit{bursts} de SSH, \textit{keepalives}
HTTP e fluxos de curta duração sobre rede local rápida. Por isso, modelos
treinados exclusivamente nos datasets sintéticos tendem a classificar
erroneamente o tráfego benigno real como ataque. Misturar esse tráfego real à
classe benigna reduz o desvio de distribuição entre treino e implantação.

Dada a escala dos conjuntos de dados modernos de IDS, que reúnem milhões de
registros e capturas brutas da ordem de terabytes, manter os dados
integralmente em memória é inviável. A estratégia de ingestão e armazenamento
adotada busca, portanto, escalabilidade, reprodutibilidade e execução em
hardware convencional, sem a necessidade de frameworks de computação
distribuída. Cada conjunto de dados é consolidado em uma única tabela
SQLite\footnote{\url{https://www.sqlite.org/}}~\cite{sqlite}, preservando a
estrutura original de \textit{features} e rótulos. O SQLite foi selecionado
devido à sua natureza leve, portabilidade e capacidade de realizar
eficientemente agregações, contagens distintas e consultas estatísticas
diretamente em disco. A construção do banco segue duas etapas automatizadas e
reprodutíveis, cujo pipeline está disponível
publicamente\footnote{\url{https://github.com/luiz-linkezio/A-Hierarchical-Solution-based-on-MEC-and-IoT-for-Cyberattack-Detection-in-6G-Networks/blob/main/notebooks/database_creation.ipynb}}:

\begin{enumerate}
  \item \textbf{PCAP $\rightarrow$ CSV (\textit{Comma-Separated Values}) via netflower:}
    para cada pasta de categoria nos datasets (e.g.\ \texttt{benign/},
    \texttt{ddos/}, \texttt{malware/}), todos os arquivos
    \texttt{.pcap}/\texttt{.pcapng} são convertidos em fluxos via
    \texttt{convert\_pcap\_to\_csv(n\_jobs=8)}, que executa oito processos em
    paralelo particionando os pacotes por chave bidirecional determinística. O
    rótulo de cada fluxo é derivado automaticamente do nome da pasta, dispensando
    rotulagem manual e evitando inconsistências de esquema entre as fontes. Os
    fluxos de todos os PCAPs de uma mesma pasta são então concatenados em um
    único arquivo \texttt{merged\_<label>.csv}.
  \item \textbf{CSV $\rightarrow$ SQLite:} cada \texttt{merged\_*.csv} é inserido
    no banco em \textit{chunks} de 50.000 linhas, evitando carregar o arquivo
    inteiro em memória. Ao final, cria-se um índice sobre a coluna
    \texttt{label}, que restringe as consultas por classe ao subconjunto relevante
    em vez de varrer toda a tabela. Por agregar múltiplos pacotes em um único
    registro de fluxo e descartar o detalhe pacote a pacote das capturas brutas,
    a representação reduz o volume em disco em cerca de $5\times$: o banco
    resultante ocupa $\sim$190~GB, ante $\sim$1~TB das PCAPs originais.
\end{enumerate}


O \textit{netflower} é uma ferramenta de extração de fluxos de rede
desenvolvida neste projeto, publicada no Python Package Index
(PyPI)\footnote{\url{https://pypi.org/project/netflower/}} e disponível em
código aberto sob licença
MIT\footnote{\url{https://github.com/luiz-linkezio/netflower}}~\cite{netflower}.
A partir de arquivos \texttt{.pcap}/\texttt{.pcapng} ou diretamente de uma
interface de rede ao vivo, ela reconstrói fluxos bidirecionais e produz
82~\textit{features} compatíveis com o conjunto do
CICFlowMeter\footnote{\url{https://github.com/ahlashkari/CICFlowMeter}}~\cite{cicflowmeter},
organizadas em: identidade de fluxo (IPs e portas de origem/destino, protocolo,
\textit{timestamp}), estatísticas de duração e de taxa (bytes/s e pacotes/s nos
sentidos \textit{forward}, \textit{backward} e combinado), distribuições de
tamanho de pacote e de tempo entre chegadas (média, desvio-padrão, mínimo,
máximo e variância), contagens das oito \textit{flags} do protocolo da camada
de transporte TCP (FIN, SYN, RST, PSH, ACK, URG, ECE, CWR), períodos
ativo/ocioso, métricas de \textit{bulk} e de \textit{subflow} e os tamanhos
iniciais da janela TCP. Apenas fluxos IPv4 sobre TCP e UDP são extraídos; os
demais protocolos são silenciosamente ignorados.

A ferramenta foi projetada desde o início para execução em dispositivos de
borda. O \textit{parsing} dos pacotes usa o
\textit{dpkt}\footnote{\url{https://github.com/kbandla/dpkt}}~\cite{dpkt},
biblioteca 10--25$\times$ mais rápida que o
Scapy\footnote{\url{https://scapy.net}}~\cite{scapy} em processadores ARM. Para
manter a memória constante por fluxo, as estatísticas são acumuladas pelo
algoritmo \textit{online} de Welford~\cite{welford}, que atualiza média e
variância em $O(1)$ sem armazenar listas de pacotes, enquanto uma coleta
periódica de fluxos ociosos (\texttt{gc\_interval}) limita o estado residente
mesmo sob grande número de fluxos simultâneos. A captura ao vivo acessa o
libpcap\footnote{\url{https://www.tcpdump.org}}~\cite{libpcap} via
\texttt{ctypes}, sem dependências de compilação no nó de borda. Por fim, cada fluxo é emitido apenas ao completar, seja pelo término da
conexão (TCP FIN/RST), seja por \textit{timeout} de inatividade
(\texttt{idle\_timeout}, padrão 30~s) ou de duração máxima
(\texttt{flow\_timeout}, padrão 120~s), garantindo \textit{features} completas
em toda linha; esses dois \textit{timeouts} são justamente os parâmetros que
governam a latência de detecção do IDS na captura ao vivo
(Seção~\ref{subsec:validacao}).

\subsection{Mapeamento de rótulos}
\label{subsec:rotulos}

Uma vez que os ataques DDoS e DoS visam exaurir recursos e seus fluxos de rede
(alta taxa de pacotes, curta duração, anomalias de \textit{flags}) são
estatisticamente similares, manter duas classes separadas pode adicionar ruído
sem ganho discriminativo significativo no classificador. Nesse sentido, após
análise das distribuições das duas classes, realizou-se a unificação delas,
reduzindo o espaço de classes total de 9 para 8 (7 de ataque + 1 benigno),
simplificando o problema da Fase~2. Para isso, aplicou-se
$\texttt{LABEL\_MAP} = \{\texttt{ddos} \rightarrow \texttt{dos}\}$ antes do
treinamento.

Do ponto de vista do nó alvo (VIM~4), a distinção entre DoS e DDoS é de natureza
\textit{externa} aos fluxos individuais: um ataque DoS provém de uma única fonte,
enquanto o DDoS usa múltiplas fontes distribuídas, mas os fluxos individuais
recebidos pelo alvo são estatisticamente indistinguíveis. A característica que
capturaria essa diferença (diversidade de IPs de origem) exige agregação
\textit{entre} fluxos e não está presente nas 82~\textit{features} por fluxo do
\textit{netflower}. Manter as duas classes separadas no treinamento produziria
um classificador que aprenderia a distingui-las por artefatos dos datasets, não
por características do tráfego real observado pelo nó de borda. A
Tabela~\ref{tab:taxonomy} apresenta a taxonomia de fluxos resultante.

\begin{table}[ht]
\centering
\caption{Taxonomia de ataques utilizada no treinamento do IDS.}
\label{tab:taxonomy}
\begin{tabular}{@{}ll@{}}
\toprule
\textbf{Categoria} & \textbf{Descrição} \\ \midrule
\texttt{benign}     & Tráfego legítimo de rede \\
\texttt{dos}        & DoS e DDoS unificados (exaustão de recursos) \\
\texttt{malware}    & Malware e comunicação de Comando e Controle (C2) \\
\texttt{bruteforce} & Ataques de força bruta (SSH, HTTP, Telnet) \\
\texttt{mitm}       & Ataques MITM (envenenamento de ARP, \textit{Address Resolution Protocol}) \\
\texttt{web}        & Ataques a aplicações web (scanning, fuzzing) \\
\texttt{spoofing}   & Falsificação de identidade de origem em pacotes IP \\
\texttt{recon}      & Varreduras e reconhecimento de rede (port scan) \\ \bottomrule
\end{tabular}
\end{table}

\subsection{Engenharia de características}
\label{subsec:features}

Partindo de 83 colunas brutas produzidas pelo \textit{netflower}, a etapa de
limpeza removeu grupos de \textit{features} que compartilham um mesmo risco:
induzir o modelo a aprender algo diferente do padrão estatístico do tráfego.

O primeiro grupo é o de identidade de fluxo (\texttt{src\_ip},
\texttt{dst\_ip}, \texttt{src\_port}, \texttt{dst\_port} e
\texttt{timestamp}). Colunas que identificam hosts e serviços específicos fazem
o modelo memorizar endereços vistos no treino em vez de aprender padrões de
tráfego, de modo que qualquer host novo no ambiente de implantação produziria
viés de previsão; todas foram removidas. Um problema análogo aparece nas colunas
de janela TCP inicial (\texttt{init\_fwd\_win\_byts} e
\texttt{init\_bwd\_win\_byts}): o \textit{netflower} escreve o valor-sentinela
$-1$ quando o tamanho da janela do \textit{handshake} não está disponível, valor
que não ocorre no tráfego real. Mantida a \textit{feature}, o modelo aprenderia
a disparar sobre o sentinela, gerando falsos positivos em fluxos legítimos nos
quais o campo é genuinamente ausente, e por isso as duas colunas foram
descartadas.

As \textit{features} de \textit{bulk rate} (\texttt{fwd\_byts\_b\_avg} e
similares) sofrem de um desvio sistemático entre treino e implantação. Elas
medem a taxa de transferência em segmentos de tráfego contínuo, o que exige que
o extrator identifique janelas de envio ininterrupto; o \textit{netflower}, em
modo ao vivo, não implementa essa heurística e emite zero para todos os fluxos,
enquanto os datasets de treino contêm valores não nulos gerados por ferramentas
que a aplicam. Manter essas colunas tornaria o modelo dependente de um sinal que
nunca existirá em produção, tornando-as não apenas inúteis, mas potencialmente
enganosas; foram, portanto, eliminadas.

Por fim, dois filtros estatísticos completaram a limpeza. Onze
\textit{features} foram removidas por correlação alta com outra ($|\rho| >
0{,}95$, Pearson): em pares altamente correlacionados, ambas carregam
essencialmente a mesma informação, e manter as duas inflaria o peso do sinal
subjacente e o custo de inferência sem ganho discriminativo; o limiar
conservador preserva \textit{features} com sobreposição apenas parcial, que
ainda podem capturar variações independentes do tráfego. Outras três foram
descartadas por variância próxima de zero (\textit{VarianceThreshold},
$\sigma^2 < 10^{-4}$): derivavam de contadores que o \textit{netflower} emite
com valor fixo na maioria dos protocolos capturados e, ao assumirem praticamente
o mesmo valor em todos os fluxos, não contribuem para separar classes, atuando
apenas como ruído constante.

O resultado foram 55 \textit{features} numéricas finais. As duas de maior
relevância pelo ganho do LightGBM foram a taxa de bytes por segundo
(\texttt{flow\_byts\_s}) e o tempo mínimo de chegada entre pacotes
(\texttt{flow\_iat\_min}), que capturam, respectivamente, a intensidade e o
padrão temporal do tráfego.

\subsection{Treinamento das duas fases}
\label{subsec:treinamento}

O treinamento ocorre sobre milhões de fluxos descritos por dezenas de
\textit{features}, mas a inferência precisa caber no orçamento computacional de
um nó de borda ARM, o que pede um algoritmo que combine acurácia em dados
tabulares de alta dimensionalidade com modelos compactos e rápidos na predição.
O LightGBM~\cite{ke:17} atende a esse perfil: é um método de \textit{gradient
boosting} sobre árvores de decisão com crescimento folha a folha
(\textit{leaf-wise}), que converge rapidamente e produz modelos leves o
suficiente para inferência em dispositivos restritos. Resta o ajuste de
hiperparâmetros, cuja busca exaustiva seria proibitiva no volume de dados
utilizado; o Optuna~\cite{akiba:19} trata esse problema com otimização baseada
em TPE (\textit{Tree-structured Parzen Estimator}) e poda antecipada de
tentativas pouco promissoras, reduzindo o custo de encontrar boas configurações.

As duas fases do IDS hierárquico, o classificador binário da Fase~1 e o
multiclasse da Fase~2, foram treinadas de forma independente sobre o banco
SQLite, mas com objetivos de otimização deliberadamente distintos, coerentes com
o papel de cada uma na arquitetura hierárquica. A Fase~1, como filtro de
primeira linha, é ajustada para não deixar ataques passarem, priorizando a
revocação por meio do escore F2. A Fase~2, acionada apenas sobre fluxos já
sinalizados, é ajustada para distribuir corretamente o tipo de ataque entre as
classes, priorizando o macro F1. Em ambos os casos, os hiperparâmetros foram
buscados automaticamente com o Optuna e a amostragem partiu diretamente do
SQLite, garantindo balanceamento entre classes. O treinamento foi executado em
CPU no PC descrito na Seção~\ref{subsec:validacao}, e as próximas subseções
descrevem a configuração e os resultados \textit{offline} de cada fase.

\subsubsection{Fase 1: classificador binário}

O classificador binário da Fase~1 foi treinado sobre 615.317 fluxos
benignos e os 554.795 fluxos de ataque disponíveis; após a remoção de fluxos
duplicados, restaram 703.558 registros, com divisão estratificada de 80/20 entre
treino e teste. Tanto os hiperparâmetros quanto o limiar de decisão foram
otimizados conjuntamente pelo Optuna (20 tentativas, 1.800~s de
\textit{timeout}), tendo o escore F2 como objetivo. A escolha do F2 é
deliberada, pois ele pondera a revocação quatro vezes mais que a precisão,
refletindo que, em um IDS (e não em um Sistema de Prevenção de Intrusão, IPS),
um falso negativo (um ataque não detectado) é mais custoso que um falso positivo
(um alarme indevido).

A Tabela~\ref{tab:phase1_offline} apresenta a avaliação no conjunto de teste
(140.712 fluxos) ao limiar ótimo de 0,157. O resultado evidencia o efeito desse
objetivo: a revocação de 96\% na classe de ataque significa que apenas 4\% dos
ataques atravessariam o filtro de primeira linha sem gerar alerta, exatamente o
comportamento esperado da fase cuja função é não deixar ataques passarem. O
custo é uma precisão de 68\%, de modo que cerca de um em cada três alertas seria
um alarme falso; em operação contínua no nó de borda, esse volume de alarmes
sobrecarregaria tanto o operador quanto a Fase~2, acionada para cada fluxo
sinalizado. O escore F2 resultante é de 88,6\%. Esse compromisso é aceitável em
treino, mas elevado demais para produção; por isso, o limiar de implantação foi
posteriormente fixado em 0,9, trocando parte da revocação por uma precisão
substancialmente maior e um volume de falsos positivos próximo de zero
(Seção~\ref{sec:resultados}).

\begin{table}[ht]
\centering
\caption{Avaliação \textit{offline} da Fase~1 no conjunto de teste (limiar ótimo = 0,157).}
\label{tab:phase1_offline}
\begin{tabular}{@{}lrrrr@{}}
\toprule
\textbf{Classe} & \textbf{Precisão} & \textbf{Revocação} & \textbf{F1} & \textbf{Suporte} \\ \midrule
attack  & 68\% & \textbf{96\%} & 80\% & 57.398 \\
benign  & \textbf{96\%} & 69\% & 80\% & 83.314 \\ \midrule
\textbf{Acurácia} & & & \textbf{80\%} & 140.712 \\
Macro avg & 82\% & 83\% & 80\% & \\
\bottomrule
\end{tabular}
\end{table}

\subsubsection{Fase 2: classificador multiclasse}

O classificador multiclasse da Fase~2 foi treinado a partir de uma
amostragem independente de até 615.317 linhas por classe extraídas diretamente do
SQLite. Para compensar a escassez da classe \textit{bruteforce}
($\sim$16~mil amostras), aplicou-se \texttt{class\_weight=balanced}, e o objetivo
de otimização do Optuna foi o macro F1, que dá peso igual a todas as classes e,
portanto, penaliza explicitamente o mau desempenho nas classes minoritárias.

A Tabela~\ref{tab:phase2_offline} reporta o desempenho por classe no conjunto de
teste (317.381 fluxos). Já em avaliação \textit{offline} emerge o padrão que se
acentuará na validação no \textit{testbed}: a classe \texttt{dos} é classificada
de forma quase perfeita (F1 = 98\%), enquanto as classes de fluxo curto e baixa
amostragem \texttt{bruteforce}, \texttt{spoofing} e \texttt{mitm} apresentam
precisão baixa (20\%, 28\% e 42\%), revelando confusão sistemática com as
demais. A revocação relativamente alta dessas classes minoritárias (72\% a
77\%) é efeito direto do \texttt{class\_weight=balanced}, que força o modelo a
não ignorá-las, mas o macro F1 de 67\% confirma que a separação entre os tipos
de ataque é desigual: o sinal por fluxo distingue bem ataques de assinatura
peculiar e mal os de comportamento semelhante a inundação. Na prática, isso
significa que o rótulo de tipo emitido pela Fase~2 é confiável para varreduras
e inundações, mas ainda não para força bruta, \textit{spoofing} e MITM, cuja
triagem segue dependendo do alerta binário da Fase~1.

\begin{table}[ht]
\centering
\caption{Avaliação \textit{offline} da Fase~2 por classe no conjunto de teste.}
\label{tab:phase2_offline}
\begin{tabular}{@{}lrrrr@{}}
\toprule
\textbf{Classe} & \textbf{Precisão} & \textbf{Revocação} & \textbf{F1} & \textbf{Suporte} \\ \midrule
benign      & 83\% & 69\% & 76\% & 83.313 \\
bruteforce  & 20\% & 77\% & 31\% &  2.191 \\
dos         & \textbf{99\%} & \textbf{96\%} & \textbf{98\%} & 64.582 \\
malware     & 90\% & 78\% & 84\% & 81.338 \\
mitm        & 42\% & 72\% & 53\% & 13.823 \\
recon       & 83\% & 68\% & 75\% & 47.915 \\
spoofing    & 28\% & 73\% & 41\% & 12.289 \\
web         & 77\% & 84\% & 80\% & 11.930 \\ \midrule
\textbf{Acurácia} & & & \textbf{78\%} & 317.381 \\
Macro F1    & & & \textbf{67\%} & \\
\bottomrule
\end{tabular}
\end{table}

\subsubsection{Exploração preliminar de uma terceira fase}
\label{subsubsec:fase3}

A concepção original da arquitetura previa uma terceira fase, de agrupamento
não supervisionado, para os fluxos que a Fase~2 classifica com baixa confiança,
candidatos a ameaças desconhecidas. Essa fase foi explorada apenas
\textit{offline}, sobre o conjunto de teste da Fase~2: os 42.527 fluxos (13,4\%
do total) cuja confiança máxima ficou abaixo de 0,4 foram projetados pelo
UMAP~\cite{mcinnes:18} em 10 dimensões (\texttt{n\_neighbors} = 30,
\texttt{min\_dist} = 0) e agrupados pelo HDBSCAN~\cite{mcinnes:17}
(\texttt{min\_cluster\_size} = 100, \texttt{min\_samples} = 50), que encontrou
84 grupos e marcou 4,7\% dos fluxos como ruído. Os modelos resultantes estão
versionados no repositório do projeto, mas a fase não foi integrada ao IDS
implantado no nó de borda nem avaliada no \textit{testbed}; sua caracterização
(a que ataques os grupos correspondem e se de fato isolam tráfego inédito)
permanece como trabalho futuro.

\subsection{Validação no \textit{testbed}: ambiente, ferramentas e protocolo}
\label{subsec:validacao}

A validação no \textit{testbed} envolveu dois equipamentos conectados na
mesma rede local, ilustrados na Figura~\ref{fig:ambiente} e especificados
na Tabela~\ref{tab:hardware}. O computador pessoal (PC) acumulou três papéis:
treinou os modelos, orquestrou todo o experimento por SSH e gerou os ataques. A
placa VIM~4 atuou como nó de borda MEC, executando os scripts IDS e monitorando a
interface de rede \texttt{eth0}, alvo de todos os ataques. A rede local Gigabit, compartilhada com outros hosts
(roteador e dispositivos domésticos), faz as vezes do enlace de acesso do cenário MEC alvo; como o IDS decide
sobre fluxos IP reconstruídos acima da camada de acesso
(Seção~\ref{subsec:pipeline}), essa substituição não altera a representação de
tráfego vista pelo detector. O treinamento do LightGBM foi executado em CPU
(8 \textit{threads}), e os modelos treinados foram copiados para a VIM~4 para
inferência.

A VIM~4\footnote{\url{https://www.khadas.com/vim4}} foi escolhida como nó de
borda por reunir, em um computador de placa única de baixo custo, os atributos
esperados de um nó MEC de acesso. Seu \textit{System-on-Chip} (SoC) Amlogic
A311D2 combina oito núcleos ARM em arranjo heterogêneo
(Tabela~\ref{tab:hardware}): os quatro Cortex-A73 a 2,2~GHz absorvem a carga
sustentada de captura e inferência, enquanto os quatro Cortex-A53 a 1,8~GHz
atendem às tarefas leves com consumo reduzido, mantendo a placa em um envelope
de potência de 2,5 a 9~W, uma fração do consumo de um servidor convencional e
compatível com implantação em pontos de acesso com alimentação restrita. Os
8~GB de RAM LPDDR4x dão margem para os modelos das duas fases residirem em
memória simultaneamente.

\begin{figure}[ht]
\centering
\definecolor{cardfill}{HTML}{F5F7FA}
\definecolor{cardstroke}{HTML}{334155}
\definecolor{devfill}{HTML}{E2E8F0}
\definecolor{devstroke}{HTML}{334155}
\definecolor{screenfill}{HTML}{93C5FD}
\definecolor{boardfill}{HTML}{CBD5E1}
\definecolor{chipfill}{HTML}{475569}
\definecolor{netblue}{HTML}{2563EB}
\definecolor{netteal}{HTML}{0F766E}
\resizebox{\textwidth}{!}{%
\begin{tikzpicture}[
  font=\footnotesize, line join=round,
  card/.style={rectangle, rounded corners=5pt, draw=cardstroke, thick,
               fill=cardfill, inner sep=9pt},
  subbox/.style={rectangle, rounded corners=3pt, draw=devstroke, semithick,
                 fill=white, align=center, font=\scriptsize, inner sep=5pt,
                 text width=2.7cm},
  logbox/.style={subbox, fill=blue!6, draw=netblue!70},
  callarr/.style={-{Stealth[length=2mm]}, thick, draw=devstroke},
  dataarr/.style={-{Stealth[length=2mm]}, thick, draw=netteal},
  attackarr/.style={-{Stealth[length=3mm]}, line width=1pt, draw=netblue,
                    dash pattern=on 0.5pt off 2.5pt},
  collectarr/.style={-{Stealth[length=2.5mm]}, line width=0.9pt, draw=netteal,
                     dash pattern=on 0.5pt off 2.5pt}]

  %======== PC: hardware + elementos de software ========
  \node (pcimg) {%
    \begin{tikzpicture}[line join=round]
      % gabinete
      \draw[thick, draw=devstroke, fill=devfill, rounded corners=1pt]
            (-0.80,-0.35) rectangle (-0.50,0.62);
      \draw[devstroke] (-0.75,0.50)--(-0.55,0.50);
      \draw[devstroke] (-0.75,0.42)--(-0.55,0.42);
      \fill[netblue] (-0.65,-0.16) circle(0.035);
      % monitor
      \draw[thick, draw=devstroke, fill=devfill, rounded corners=2pt]
            (-0.35,0.00) rectangle (0.80,0.62);
      \fill[screenfill] (-0.28,0.08) rectangle (0.73,0.54);
      \draw[thick, devstroke] (0.22,0.00)--(0.22,-0.18);
      \draw[line width=1.4pt, devstroke] (0.02,-0.20)--(0.42,-0.20);
    \end{tikzpicture}};
  \node[below=3pt of pcimg, align=center] (pctitle)
    {\textbf{PC}\\Ryzen~5 7600 + RTX~5070\\{\scriptsize treino dos modelos (CPU)}};
  \node[subbox, below=8pt of pctitle] (orch)
    {\textbf{Orquestrador}\\\texttt{\scriptsize run\_experiment.sh}\\[3pt]
     \tiny deploy SSH, sessões A/B, teardown, métricas};
  \node[subbox, below=9pt of orch] (atk)
    {\textbf{Gerador de ataques}\\\texttt{\scriptsize attack\_generator.py}\\[3pt]
     \tiny 8 categorias de ataque};
  \node[logbox, below=9pt of atk] (relat)
    {\textbf{Relatório de ataques}\\\texttt{\scriptsize report\_*.json}\\[3pt]
     \tiny janelas de \textit{ground truth} por ataque};
  \node[subbox, below=9pt of relat] (metrics)
    {\textbf{Métricas}\\\texttt{\scriptsize ids\_metrics.py}\\[3pt]
     \tiny cruza log do IDS + relatório; acurácia, F1, energia};
  \draw[callarr] (orch) -- node[right, font=\tiny, text=black, xshift=1pt]{invoca} (atk);
  \draw[dataarr] (atk) -- node[right, font=\tiny, text=black, xshift=1pt]{gera} (relat);
  \draw[dataarr] (relat) -- node[right, font=\tiny, text=black, xshift=1pt]{coleta} (metrics);

  %======== VIM 4: hardware + elementos de software ========
  \node[right=7.5cm of pcimg] (vimimg) {%
    \begin{tikzpicture}[line join=round]
      % placa (PCB)
      \draw[thick, draw=devstroke, fill=boardfill, rounded corners=2pt]
            (-0.82,-0.52) rectangle (0.82,0.52);
      \foreach \x/\y in {-0.70/0.40, 0.70/0.40, -0.70/-0.40, 0.70/-0.40}
        \draw[devstroke, fill=white] (\x,\y) circle(0.05);
      % SoC + dissipador
      \draw[thick, draw=devstroke, fill=chipfill] (-0.28,-0.28) rectangle (0.22,0.28);
      \foreach \x in {-0.22,-0.11,0.00,0.11,0.18}
        \draw[white, line width=0.6pt] (\x,-0.24)--(\x,0.24);
      % barra de pinos (GPIO)
      \foreach \x in {-0.58,-0.49,-0.40,-0.31}
        \fill[chipfill] (\x,0.36) rectangle ++(0.05,0.10);
      % portas (USB / Ethernet)
      \draw[thick, draw=devstroke, fill=devfill] (0.52,0.12) rectangle (0.86,0.34);
      \draw[thick, draw=devstroke, fill=devfill] (0.52,-0.16) rectangle (0.86,0.06);
    \end{tikzpicture}};
  \node[below=3pt of vimimg, align=center] (vimtitle)
    {\textbf{VIM~4}\\Amlogic A311D2 (ARM)\\{\scriptsize nó de borda MEC}};
  \node[subbox, below=8pt of vimtitle] (idsbox)
    {\textbf{IDS}\\\texttt{\scriptsize network\_ids.py}\\[3pt]
     \tiny Fase~1 (+ Fase~2)\\interface \texttt{eth0}};
  \node[logbox, below=9pt of idsbox] (idslog)
    {\textbf{Log do IDS}\\\texttt{\scriptsize ids\_run\_*.log}\\[3pt]
     \tiny alertas por fluxo + \textit{snapshots} de sistema};
  \draw[dataarr] (idsbox) -- node[right, font=\tiny, text=black, xshift=1pt]{registra} (idslog);

  %======== cartões de fundo ========
  \begin{scope}[on background layer]
    \node[card, fit=(pcimg)(pctitle)(orch)(atk)(relat)(metrics)] (pccard) {};
    \node[card, fit=(vimimg)(vimtitle)(idsbox)(idslog)] (vimcard) {};
  \end{scope}

  %======== switch / rede isolada (centro) ========
  \coordinate (midx) at ($(pccard.east)!0.5!(vimcard.west)$);
  \node (sw) at (midx |- idsbox) {%
    \begin{tikzpicture}[line join=round]
      \draw[thick, draw=devstroke, fill=devfill, rounded corners=1.5pt]
            (-0.52,-0.16) rectangle (0.52,0.16);
      \foreach \x in {-0.40,-0.24,-0.08,0.08,0.24}
        \draw[devstroke, fill=white] (\x,-0.10) rectangle ++(0.10,0.12);
      \fill[netblue] (0.45,0.09) circle(0.022);
    \end{tikzpicture}};
  \node[above=1pt of sw, font=\scriptsize\itshape, text=black]
    {Rede local (Gigabit Ethernet)};

  %======== conexões entre os nós (rede isolada) ========
  % tráfego de ataque: sai do gerador de ataques direto para o IDS da VIM
  \draw[attackarr] (atk.east) --
        node[pos=0.5, below=3pt, align=center, font=\scriptsize, text=black]
        {ataques (8 categorias)\\\texttt{192.168.100.13}}
        (idsbox.west);
  % coleta de log: ids\_metrics.py busca o log do IDS na VIM 4 via SSH
  \draw[collectarr] (idslog.west) -- ++(-0.7,0) |- (metrics.east);
  \node[above=2pt, font=\scriptsize, text=black] at (midx |- metrics)
    {coleta do log (via SSH)};
\end{tikzpicture}}
\caption{Ambiente de experimentação: PC e VIM~4 conectados pela mesma rede local.}
\label{fig:ambiente}
\end{figure}

\begin{table}[ht]
\centering
\caption{Especificações do hardware utilizado.}
\label{tab:hardware}
\begin{tabular}{@{}lll@{}}
\toprule
\textbf{Componente} & \textbf{PC (treino / ataques)} & \textbf{VIM 4 (IDS)} \\ \midrule
SoC / CPU & AMD Ryzen 5 7600 (6c/12t) & Amlogic A311D2 \\
          &                            & 4$\times$ Cortex-A73 @ 2,2~GHz \\
          &                            & + 4$\times$ Cortex-A53 @ 1,8~GHz \\
GPU       & NVIDIA RTX 5070 (16~GB)    & Mali-G52 MP8 \\
RAM       & 2$\times$16~GB DDR5 5600~MHz & 8~GB LPDDR4x \\
Placa-mãe & Asus TUF Gaming B650 Plus WiFi & --- \\
Fonte     & MSI A750GF (750~W, 80+ Gold) & --- \\
Rede      & Gigabit Ethernet           & Gigabit Ethernet \\
\bottomrule
\end{tabular}
\end{table}

A Figura~\ref{fig:ambiente} detalha os elementos de software que executam sobre
esse hardware. No PC residem quatro elementos. O orquestrador
(\texttt{run\_experiment.sh}) coordena todo o experimento: implanta os scripts na
VIM~4 via SSH, conduz as duas sessões de avaliação (A e B, descritas adiante),
encerra os processos ao final de cada sessão e dispara o cálculo de métricas. O
gerador de ataques (\texttt{attack\_generator.py}), invocado pelo orquestrador,
executa as oito categorias de ataque em sequência contra a VIM~4. Ao término de
cada execução, o gerador produz o relatório de ataques (\texttt{report\_*.json}),
que registra as janelas temporais de início e fim de cada ataque e constitui o
\textit{ground truth} da avaliação. Por fim, o script de métricas
(\texttt{ids\_metrics.py}) coleta via SSH o log produzido pelo IDS na VIM~4 e o
cruza com o relatório de ataques para calcular as métricas de detecção, de
recursos e de energia. Na VIM~4, por sua vez, executa o IDS propriamente dito
(\texttt{network\_ids.py}), que monitora a interface \texttt{eth0}, roda a Fase~1
(e, quando habilitada, a Fase~2) e registra sua operação no log do IDS
(\texttt{ids\_run\_*.log}), contendo os alertas emitidos por fluxo e
\textit{snapshots} periódicos do estado do sistema (CPU, RAM, temperatura e
frequência), dos quais derivam as métricas de custo. Para não estagnar a captura
sob \textit{flood}, o IDS separa a captura da inferência: a \textit{thread} de
captura apenas enfileira cada fluxo concluído, e uma \textit{thread} distinta
consome a fila em lote, rodando a Fase~1 sobre o lote e a Fase~2 apenas sobre os
fluxos alertados. Sem essa separação, a inferência bloquearia a captura e o
\textit{buffer} do \textit{kernel} transbordaria, descartando pacotes. A fila não
tem limite prático, de modo que todo fluxo capturado é classificado ainda que o
processamento fique para trás do ataque, e sua profundidade é registrada no log.

Para executar as 8 categorias de ataque, o gerador emprega as seguintes
ferramentas:

\begin{itemize}
  \item \textbf{nmap}\footnote{\url{https://nmap.org}}~\cite{nmap}: varredura de
    rede: SYN scan (\texttt{-sS}), UDP scan dos 100 top ports (\texttt{-sU}),
    varredura agressiva (\texttt{-A});
  \item \textbf{masscan}\footnote{\url{https://github.com/robertdavidgraham/masscan}}~\cite{masscan}:
    varredura SYN de alta velocidade;
  \item \textbf{hping3}\footnote{\url{http://www.hping.org}}~\cite{hping3}:
    geração de inundação: SYN flood (porta 80), UDP flood (porta 53), ICMP flood;
    DDoS com fonte aleatória (\texttt{--rand-source}) nas portas 443 e 80;
    spoofing com fonte forjada fixa (\texttt{-a});
  \item \textbf{medusa}\footnote{\url{https://github.com/jmk-foofus/medusa}}~\cite{medusa}:
    força bruta de credenciais: SSH (root), HTTP (admin), Telnet (root);
  \item \textbf{nikto}\footnote{\url{https://github.com/sullo/nikto}}~\cite{nikto}:
    varredura de vulnerabilidades web;
  \item \textbf{gobuster}\footnote{\url{https://github.com/OJ/gobuster}}~\cite{gobuster}:
    enumeração de diretórios/caminhos HTTP;
  \item \textbf{curl}\footnote{\url{https://curl.se}}~\cite{curl}: inundação HTTP
    (100 requisições em paralelo);
  \item \textbf{arpspoof} (suíte \textit{dsniff})\footnote{\url{https://www.monkey.org/~dugsong/dsniff/}}~\cite{dsniff}:
    envenenamento ARP bidirecional (ataque MITM), com tráfego forçado pela rota
    envenenada;
  \item \textbf{nc} (netcat)\footnote{\url{https://nmap.org/ncat/}}~\cite{ncat}:
    beacon C2 simulado via TCP porta~4444 (100 conexões curtas).
\end{itemize}

Os dois pipelines são avaliados em duas sessões sequenciais sob o mesmo
roteiro: a Sessão~A apenas com o IDS binário e a Sessão~B com o
binário+multiclasse, o que dá atribuição limpa de recursos a cada pipeline e
janelas de \textit{ground truth} estruturalmente idênticas. Cada sessão inclui
60~s de baseline benigno antes e depois dos ataques e 120~s de intervalo ocioso
entre ataques consecutivos, que mantém as janelas de ataque separadas. A
rotulagem pelo início do fluxo não depende desse intervalo, o que importa porque
o atraso de emissão medido supera os 120~s (Seção~\ref{subsec:res_artefatos}).
Os ataques visam
\texttt{192.168.100.13} (VIM~4), com sincronização por \textit{offset} de 3~h de
BRT para UTC; limiar de detecção binária $P_1 \geq 0{,}9$; \textit{netflower}
com \texttt{idle\_timeout} = 30~s e \texttt{flow\_timeout} = 120~s. A captura
O \textit{netflower} converte apenas tráfego TCP e UDP em fluxos, de modo que
os \textit{floods} ICMP de DoS e \textit{spoofing} não geram fluxos e as
contagens desses ataques cobrem só a parte TCP/UDP. A captura
PCAP no atacante foi deliberadamente desativada: as métricas derivam do log do
IDS (todos os fluxos classificados + amostras periódicas de sistema) e das janelas temporais do
relatório do gerador de ataques, não dos PCAPs, que, sob os \textit{floods},
chegavam a vários GB sem valor analítico.

A detecção é avaliada \emph{por fluxo}: cada fluxo classificado pelo IDS é
rotulado como ataque ou benigno segundo a janela de \textit{ground truth} em que
\emph{começou} (\texttt{flow\_ts}, o instante de captura do primeiro pacote no
cabeçalho do PCAP), e nunca segundo o instante em que o alerta foi emitido. A
escolha é deliberada: sob \textit{flood}, o nó ARM não classifica em tempo real e
os fluxos são emitidos com atraso mediano de 180--437~s
(Seção~\ref{subsec:res_artefatos}), de modo que qualquer métrica ancorada na
emissão, inclusive a contagem por segundo, contamina segundos e janelas vizinhos. Rotulando pelo início do fluxo, o atraso
torna-se irrelevante para a detecção. Um filtro de \textit{endpoints} distingue
o tráfego de fundo (multicast, outros \textit{hosts} da LAN, serviços da própria
VIM) do tráfego de ataque, para que fluxos benignos que caem numa janela contem
como benignos. Reportam-se duas definições da janela: a \textbf{estrita}, com os
instantes registrados pelo orquestrador, e a \textbf{calibrada} ($+2$~s), que os
alarga em 2~s de cada lado para corrigir a defasagem de registro do orquestrador
(que grava o \texttt{start\_time} depois de já ter disparado a ferramenta, de
modo que os primeiros fluxos de ataque começam até $\sim$2~s antes do instante
anotado). A calibração continua rotulando pelo \texttt{flow\_ts}. A avaliação
multiclasse prevê o tipo da Fase~2 quando a Fase~1 alerta, e benigno caso
contrário. CPU, RAM e \textit{throughput} derivam das amostras periódicas do
log; a energia segue um modelo linear de CPU calibrado, detalhado na
Seção~\ref{subsec:res_comp}.

\subsection{Métricas de avaliação}
\label{subsec:metricas}

Convém precisar a semântica das métricas reportadas e por que são adequadas a
um IDS. Nas equações a seguir, TP, FP, TN e FN denotam, respectivamente, os
verdadeiros positivos, falsos positivos, verdadeiros negativos e falsos
negativos. A precisão é a fração de alertas que correspondem a ataques reais,
\begin{equation}
\text{Precisão} = \frac{\mathrm{TP}}{\mathrm{TP} + \mathrm{FP}},
\label{eq:precisao}
\end{equation}
de modo que uma precisão baixa significa muitos alarmes falsos, que
sobrecarregam o operador. A revocação (ou taxa de verdadeiros positivos, TPR) é
a fração de ataques efetivamente detectados,
\begin{equation}
\text{Revocação} = \frac{\mathrm{TP}}{\mathrm{TP} + \mathrm{FN}},
\label{eq:revocacao}
\end{equation}
e uma revocação baixa significa ataques que passam despercebidos. As duas se
complementam, e o F1 é sua média harmônica, um resumo único que só é alto quando
ambas o são:
\begin{equation}
\mathrm{F1} = 2 \times \frac{\text{Precisão} \times \text{Revocação}}{\text{Precisão} + \text{Revocação}}.
\label{eq:f1}
\end{equation}
A taxa de falsos positivos (FPR) é a fração de tráfego benigno classificado
como ataque,
\begin{equation}
\mathrm{FPR} = \frac{\mathrm{FP}}{\mathrm{FP} + \mathrm{TN}},
\label{eq:fpr}
\end{equation}
que em um IDS é crítica, pois um FPR alto inviabiliza a operação ao afogar a
equipe em alarmes sobre tráfego legítimo. A taxa de falsos negativos (FNR) mede
o complemento, os ataques perdidos:
\begin{equation}
\mathrm{FNR} = \frac{\mathrm{FN}}{\mathrm{FN} + \mathrm{TP}} = 1 - \text{Revocação}.
\label{eq:fnr}
\end{equation}
Como num IDS um ataque não detectado é mais custoso que um alarme indevido, este
trabalho privilegia revocação e baixo FPR, e usa F1 e macro F1 (média do F1 por
classe, que dá peso igual a classes raras) para avaliar o equilíbrio entre os
tipos de ataque na Fase~2. O escore F2 usado no treinamento da Fase~1 é a
generalização do F1 que pondera a revocação quatro vezes mais que a precisão.

% ============================================================================
\section{APRESENTAÇÃO E DISCUSSÃO OBJETIVA DOS RESULTADOS OBTIDOS}
\label{sec:resultados}

Esta seção apresenta os resultados da validação no \textit{testbed} no VIM~4,
cujo ambiente, organização do experimento e protocolo foram descritos na
Seção~\ref{subsec:validacao}, com as métricas definidas na
Seção~\ref{subsec:metricas}. A exposição segue do particular ao geral: primeiro
avaliam-se as duas fases isoladamente, o IDS binário da Fase~1
(Seção~\ref{subsec:res_bin}) e o IDS multiclasse das Fases~1+2
(Seção~\ref{subsec:res_multi}). Em seguida, caracterizam-se os artefatos de
implantação que condicionam esses números (Seção~\ref{subsec:res_artefatos}).
Por fim, confrontam-se os dois pipelines em uma comparação direta de desempenho
de detecção e consumo de recursos, \textit{throughput} e energia
(Seção~\ref{subsec:res_comp}). Todos os números derivam de uma única execução
de cada sessão (rodada \texttt{20260930\_014420}), com os modelos treinados na
Seção~\ref{subsec:treinamento} e sem retreino.

\subsection{IDS binário (Fase 1)}
\label{subsec:res_bin}

A Tabela~\ref{tab:bin_metrics} resume o IDS binário (Fase~1) sob a métrica
\emph{por fluxo}: cada fluxo classificado é rotulado como ataque ou benigno
segundo a janela de \textit{ground truth} em que \emph{começou}
(\texttt{flow\_ts}, o instante de captura do primeiro pacote no cabeçalho do
PCAP), e nunca segundo o instante em que o alerta foi emitido. A escolha é
deliberada: sob \textit{flood}, o nó ARM não acompanha o volume de pacotes e
emite os fluxos com atraso mediano de 180~s (Seção~\ref{subsec:res_artefatos}),
de modo que qualquer métrica ancorada na emissão contamina segundos e janelas
vizinhas. Rotulando pelo início do fluxo, esse atraso torna-se irrelevante para a
detecção.

Reportam-se duas definições da janela de \textit{ground truth}. A
\textbf{estrita} usa os instantes registrados pelo orquestrador. A
\textbf{calibrada} ($+2$~s) alarga cada janela em 2~s de cada lado para corrigir
uma defasagem sistemática: o orquestrador grava o \texttt{start\_time}
\emph{depois} de já ter disparado a ferramenta de ataque, de modo que os
primeiros fluxos de ataque começam até $\sim$2~s antes do instante anotado. Sob a
janela estrita, esses fluxos, todos oriundos do atacante ou de origem
forjada, caem fora da janela e são contados como falsos positivos, embora
sejam ataque real: 88\% dos ``falsos positivos'' são exatamente isso, e a folga
de $+2$~s os reabsorve por completo (a correção satura em 2~s). A calibração
continua rotulando pelo \texttt{flow\_ts}, sem reintroduzir hora de emissão.

Sob a janela calibrada, o detector alcança precisão de 98,4\% e revocação de
64,3\%, com F1 de 77,8\%. A \emph{precisão} é a métrica de especificidade
confiável neste \textit{testbed}: como a captura é dominada por \textit{flood}
($\sim$99\% dos fluxos são ataque), o tráfego genuinamente benigno soma apenas
algumas centenas de fluxos, e o FPR (33,2\% calibrado) herda um denominador
minúsculo e ruidoso, não comparável a um FPR medido sobre um \textit{baseline}
benigno limpo. A revocação de 64\% reflete o limiar operacional de 0,9: os
\textit{floods} UDP de pacote único, sem estado bidirecional, recebem confiança de
ataque em torno de 0,81, logo abaixo do limiar, e são contados como falsos
negativos. É uma perda concentrada nas classes de pacote único, não uma cegueira
geral do detector.

\begin{table}[ht]
\centering
\caption{Matriz de confusão por fluxo do IDS binário (limiar = 0,9), sob janela
estrita e calibrada ($+2$~s).}
\label{tab:bin_cm}
\begin{tabular}{@{}llrr@{}}
\toprule
& & \multicolumn{2}{c}{\textbf{Predição (estrita / \mbox{$+2$~s})}} \\
\cmidrule{3-4}
& & \textbf{Ataque} & \textbf{Benigno} \\ \midrule
\multirow{2}{*}{\textbf{Ground truth}}
  & Ataque  & 483.815 / 518.647 & 272.329 / 287.585 \\
  & Benigno &  22.769 / \ 8.281 &  21.724 / 16.641 \\ \bottomrule
\end{tabular}
\end{table}

\begin{table}[ht]
\centering
\caption{Métricas por fluxo do IDS binário no \textit{testbed} (limiar = 0,9),
sob janela de \textit{ground truth} estrita e calibrada ($+2$~s). A precisão é a
métrica de especificidade confiável. O FPR tem denominador diminuto (poucas
centenas de fluxos benignos em captura saturada de \textit{flood}).}
\label{tab:bin_metrics}
\begin{tabular}{@{}lrr@{}}
\toprule
\textbf{Métrica} & \textbf{Estrita} & \textbf{Calibrada ($+2$~s)} \\ \midrule
Acurácia            & 63,1\% & \textbf{64,4\%} \\
Precisão            & 95,5\% & \textbf{98,4\%} \\
Revocação / TPR     & 64,0\% & \textbf{64,3\%} \\
F1-score            & 76,6\% & \textbf{77,8\%} \\
FPR                 & 51,2\% & 33,2\% \\
FNR                 & 36,0\% & 35,7\% \\
\bottomrule
\end{tabular}
\end{table}

A revocação não se distribui por igual entre os ataques
(Figura~\ref{fig:detection_rate}, janela calibrada). As inundações volumétricas
de sessão TCP são detectadas de forma robusta (\textit{spoofing} com 94\% dos
fluxos e DoS com 92\%), enquanto o DDoS cai a 46\%: o \texttt{--rand-source} do
DDoS produz sobretudo fluxos UDP de pacote único, sem estado bidirecional, que o
modelo pontua logo abaixo do limiar de 0,9. O \textit{recon} atinge 57\%: o
\textit{scan} do nmap emite poucos pacotes por porta, que não acumulam
\textit{features} suficientes antes de o fluxo expirar (\texttt{idle\_timeout} de
30~s). Os ataques de baixo volume são detectados apenas em parte: força bruta 14\%
(258 fluxos), MITM 10\% (192 fluxos), \textit{web} 0,7\% (8.154 fluxos) e
\textit{malware} 3\% (103 fluxos). Nesses casos a limitação é do modelo, e não do
gerador de tráfego: o \textit{web} gera fluxos HTTP curtos que o modelo aprende
como benignos, a força bruta são poucas tentativas de autenticação por segundo, e
o MITM, após o envenenamento ARP, encaminha tráfego TCP e UDP da vítima que o
detector não distingue de tráfego legítimo de saída.

\begin{figure}[ht]
\centering
\begin{tikzpicture}
\begin{axis}[
  xbar,
  width=0.8\textwidth, height=6.0cm,
  font=\footnotesize,
  xmin=0, xmax=112,
  xtick={0,20,40,60,80,100},
  symbolic y coords={web, malware, mitm, bruteforce, ddos, recon, dos, spoofing},
  ytick=data,
  enlarge y limits=0.08,
  xlabel={Revocação por tipo de ataque (\% dos fluxos, janela $+2$~s)},
  bar width=7pt,
  nodes near coords,
  every node near coord/.append style={font=\footnotesize,
    /pgf/number format/fixed, /pgf/number format/fixed zerofill,
    /pgf/number format/precision=1, /pgf/number format/use comma},
  xmajorgrids,
  axis lines*=left,
]
\addplot[fill=gray!45, draw=black] coordinates {
  (0.7,web) (2.9,malware) (10.4,mitm) (14.0,bruteforce)
  (46.1,ddos) (57.4,recon) (92.2,dos) (94.0,spoofing)};
\end{axis}
\end{tikzpicture}
\caption{Revocação por fluxo por tipo de ataque, IDS binário, janela calibrada
$+2$~s (Sessão~A). As inundações de sessão TCP (\textit{spoofing}, DoS) são
detectadas de forma robusta; o DDoS, dominado por fluxos UDP de pacote único, e os
ataques de baixo volume (força bruta, MITM, \textit{web}, \textit{malware}) são
detectados apenas em parte no limiar de 0,9.}
\label{fig:detection_rate}
\end{figure}

\subsection{IDS multiclasse (Fases 1+2)}
\label{subsec:res_multi}

As oito classes geraram tráfego, com volumes muito desiguais: as inundações
volumétricas (DoS, DDoS e \textit{spoofing}) e o \textit{recon} dominam os fluxos,
enquanto força bruta (256 fluxos), MITM (116) e \textit{malware} (103) somam poucas
centenas. Como cada fluxo é rotulado pela janela em que começou, os resultados por
classe são atribuíveis a cada ataque, e não a alertas atrasados de ataques
anteriores.

O desempenho por classe (Tabela~\ref{tab:multi_metrics}), resumido visualmente
pela Figura~\ref{fig:multi_f1}, divide os ataques em dois regimes nítidos. A
Fase~2 classifica corretamente o \texttt{recon} (F1 = 73,7\%), enquanto as demais
classes de ataque colapsam sobre \texttt{dos}, que absorve quase todos os fluxos
alertados (F1 de 61,8\%, precisão de 74,5\% e revocação de 52,8\%). O
\texttt{spoofing} é o erro dominante: dos seus 119.873 fluxos, 98.335 recebem o
rótulo \texttt{dos} e 17.491 o rótulo \texttt{malware}, de modo que sua revocação
de tipo é praticamente nula. A acurácia de tipo é de 44,2\% (macro-F1 de 19,4\%
sobre as classes de ataque). Restrita aos fluxos que a Fase~1 de fato detectou, a
acurácia de tipo sobe para 68,6\%. Na prática, o rótulo de tipo distingue
varreduras de reconhecimento, mas não separa os demais ataques de uma inundação
genérica, o que mantém a triagem dependente do alerta binário da Fase~1 nesses
casos.

\begin{table}[ht]
\centering
\caption{IDS multiclasse na validação no \textit{testbed} (janela calibrada
$+2$~s; ddos unificado a dos).}
\label{tab:multi_metrics}
\begin{tabular}{@{}lrrrrr@{}}
\toprule
\textbf{Classe} & \textbf{Sup.} & \textbf{TP} & \textbf{FP} & \textbf{Prec.} & \textbf{F1} \\ \midrule
recon             &   1.641 &   1.005 &       83 & 92\% & \textbf{73,7\%} \\
dos (incl.\ ddos) & 625.147 & 330.203 & 113.288 & 74\% & 61,8\% \\
spoofing          & 119.873 &       0 &       28 &  0\% & 0\% \\
web               &   8.143 &      11 &        0 & 100\% & 0,3\% \\
bruteforce        &     256 &       1 &      580 &  0\% & 0,2\% \\
mitm              &     116 &       0 &    6.628 &  0\% & 0\% \\
malware           &     103 &       1 &   46.222 &  0\% & 0\% \\ \midrule
\textbf{Acurácia geral (tipo)} & \multicolumn{5}{r}{\textbf{44,2\%}} \\
\textbf{Macro F1 (classes de ataque)} & \multicolumn{5}{r}{19,4\%} \\
\bottomrule
\end{tabular}
\end{table}

\begin{figure}[ht]
\centering
\begin{tikzpicture}
\begin{axis}[
  xbar,
  width=0.8\textwidth, height=6.0cm,
  font=\footnotesize,
  xmin=0, xmax=112,
  xtick={0,20,40,60,80,100},
  xticklabel style={/pgf/number format/fixed, /pgf/number format/use comma},
  symbolic y coords={mitm, malware, spoofing, web, bruteforce, dos, recon},
  ytick=data,
  enlarge y limits=0.10,
  xlabel={F1 por tipo de ataque (\%, janela $+2$~s)},
  bar width=7pt,
  nodes near coords,
  every node near coord/.append style={font=\footnotesize,
    /pgf/number format/fixed, /pgf/number format/fixed zerofill,
    /pgf/number format/precision=1, /pgf/number format/use comma},
  xmajorgrids,
  axis lines*=left,
]
\addplot[fill=gray!45, draw=black] coordinates {
  (0.0,mitm) (0.0,malware) (0.0,spoofing) (0.3,web)
  (0.2,bruteforce) (61.8,dos) (73.7,recon)};
\end{axis}
\end{tikzpicture}
\caption{F1 por tipo de ataque do IDS multiclasse na validação no
\textit{testbed} (dados da Tabela~\ref{tab:multi_metrics}): apenas o \texttt{recon}
é bem separado; as demais classes de ataque colapsam sobre \texttt{dos}, que
absorve os fluxos alertados.}
\label{fig:multi_f1}
\end{figure}

Esse colapso tem duas causas que se somam. A primeira é a limitação do próprio
espaço de \textit{features}: no vetor de 55 atributos por fluxo, ataques de
fluxos curtos e de alta taxa tornam-se estatisticamente indistinguíveis de uma
inundação. Para \texttt{spoofing}, classificar como \texttt{dos} é inclusive a
resposta correta, já que se trata de um \textit{flood} com origem forjada.
Separar os demais exigiria estado agregado entre fluxos, como contadores de falha
de autenticação ou a diversidade de IPs de origem, ausentes por construção das
\textit{features} por fluxo. A segunda causa, para \texttt{web}, força bruta e
MITM, é o desvio entre treino e produção: na avaliação \textit{offline}
(Tabela~\ref{tab:phase2_offline}) o modelo atinge F1 de 80\%, 31\% e 53\% nessas
classes, mas ao vivo o \texttt{web} é majoritariamente classificado como benigno,
e o baixo volume da força bruta e do MITM deixa poucos fluxos para o
classificador de tipo separar.

\subsection{Desvio de distribuição e artefatos de implantação}
\label{subsec:res_artefatos}

Três fenômenos de implantação, ausentes de avaliações puramente \textit{offline},
condicionam esses números e são necessários para reproduzir o sistema. O mais
consequente é o \emph{teto de vazão} do nó de borda. A ~1.050 fluxos/s de
\textit{flood}, o VIM~4 não classifica os fluxos em tempo real. Para não descartar
tráfego, o IDS enfileira os fluxos que chegam e processa a inferência, em lote,
numa \textit{thread} separada da captura (Seção~\ref{subsec:validacao}). A fila
absorve o \textit{flood}, mas cresce ao longo de todo o ataque: o pico foi de
456.681 fluxos no experimento binário e 495.051 no multiclasse, e a latência
mediana entre o início de um fluxo e o instante em que ele é classificado subiu
para 180~s no binário e 437~s no multiclasse (máximo de 494 e 588~s). O alerta
chega, portanto, minutos depois de o fluxo ter começado. Esse deslocamento
invalida qualquer métrica ancorada no instante de emissão, inclusive a contagem
por segundo. É por isso que a detecção é avaliada \emph{por fluxo}, rotulando cada
fluxo pelo seu início (\texttt{flow\_ts}) e não pela emissão, o que a torna imune
ao atraso. O fenômeno é um limite real de capacidade do nó ARM sob ataque
volumétrico, quantificado na comparação da Seção~\ref{subsec:res_comp} e retomado
como limitação e trabalho futuro na conclusão.

O segundo é o desvio de distribuição no tráfego benigno. Os datasets
CIC utilizam tráfego benigno \textit{scripted} com padrões temporais regulares,
enquanto o tráfego interativo real (SSH e HTTP de curta duração) apresenta
características distintas que o modelo classificava como ataque. As duas
\textit{features} de maior ganho respondem pela maior parte desse erro:
\texttt{flow\_iat\_min} está na faixa de milissegundos no treino benigno, mas em
\textit{bursts} TCP reais cai para 1--10~µs, de modo que qualquer \textit{burst}
no ambiente real parece anômalo ao modelo; do mesmo modo, fluxos curtos sobre
rede local rápida produzem \texttt{flow\_byts\_s} muito acima do observado no
treino. Embora tráfego benigno real já tenha sido incorporado ao treino
(Seção~\ref{subsec:pipeline}), o desvio residual indica que a solução mais
completa é capturá-lo no próprio ambiente de implantação, a rede do VIM~4.

O terceiro fenômeno é um artefato do \textit{netflower}: cada encerramento de
conexão TCP gera um micro-fluxo adicional com $\texttt{flow\_duration} = 0$,
$\texttt{tot\_fwd\_pkts} = 1$, $\texttt{totlen\_fwd\_pkts} = 52$~bytes e apenas
as \textit{flags} FIN e ACK. Como esse padrão não tem correspondente no treino
benigno, o modelo lhe atribui sistematicamente cerca de 84\% de confiança de
ataque, uma fonte de falsos positivos que independe do tráfego real. A mitigação
é uma guarda de pacote mínimo (\texttt{tot\_fwd\_pkts} $\geq 2$) no script do
IDS, que descarta esses micro-fluxos antes da inferência.

\subsection{Comparação entre IDS binário e IDS multiclasse}
\label{subsec:res_comp}

A Tabela~\ref{tab:comparison} consolida a comparação entre os dois pipelines
testados no VIM~4. As duas sessões usam o mesmo modelo de Fase~1 e o mesmo
roteiro de ataques, mas geram tráfego de ataque próprio, por isso as métricas de
Fase~1 diferem entre elas (a Seção~\ref{subsec:res_bin} mede a Sessão~A de forma
dedicada). Adicionalmente, como a VIM~4 não expõe interfaces de
telemetria de energia em software (RAPL, \textit{Running Average Power Limit},
ou o sensor INA226), adota-se uma estimativa calibrada e não medição direta. O
método e suas limitações são descritos ao final desta seção.

\begin{table}[ht]
\centering
\caption{Comparação dos dois pipelines avaliados no VIM~4: binário (Fase~1) vs.\
hierárquico completo (Fases~1+2). As colunas são as duas sessões independentes,
que compartilham o modelo de Fase~1 mas geram tráfego próprio, por isso as
métricas de Fase~1 diferem. $^\dagger$O FPR tem denominador diminuto (poucas
centenas de fluxos benignos em captura saturada de \textit{flood}). A precisão é
a métrica de especificidade confiável.}
\label{tab:comparison}
\begin{tabular}{@{}lrr@{}}
\toprule
\textbf{Métrica} & \shortstack[r]{\textbf{IDS Binário}\\ \textbf{(só Fase~1)}}
                 & \shortstack[r]{\textbf{IDS Multiclasse}\\ \textbf{(Fases~1+2)}} \\ \midrule
\multicolumn{3}{l}{\textit{Detecção binária (Fase 1), por fluxo, janela $+2$~s}} \\
\quad Precisão            & \textbf{98,4\%} & \textbf{97,0\%} \\
\quad Revocação / TPR     & 64,3\%          & 64,0\% \\
\quad F1-score            & 77,8\%          & 77,1\% \\
\quad Acurácia            & 64,4\%          & 63,5\% \\
\quad FPR\rlap{$^\dagger$} & 33,2\%         & 47,7\% \\ \midrule
\multicolumn{3}{l}{\textit{Classificação de tipo (Fase 2)}} \\
\quad Acurácia (tipo)   & ---             & \textbf{44,2\%} \\
\quad Macro F1 (ataques) & ---            & 19,4\% \\
\midrule
\multicolumn{3}{l}{\textit{Recursos (VIM 4, durante a sessão)}} \\
\quad CPU média / máx.  & 11,8\% / 44,0\% & \textbf{40,1\% / 73,9\%} \\
\quad RAM média         & 1.907 MB        & 2.010 MB \\
\quad RAM máxima        & 2.528 MB        & 2.663 MB \\ \midrule
\multicolumn{3}{l}{\textit{Vazão, fila e latência}} \\
\quad Fluxos classificados & 831.154 & 786.821 \\
\quad Alertas emitidos  & 526.928 & 498.066 \\
\quad Tráfego de ataque & $\sim$1.078 fluxos/s & $\sim$1.021 fluxos/s \\
\quad Pico da fila      & 456.681 & 495.051 \\
\quad Latência média / máx.\ & 120 / 181~s & \textbf{352 / 550~s} \\ \midrule
\multicolumn{3}{l}{\textit{Energia estimada (modelo calibrado, Eq.~\ref{eq:power})}} \\
\quad Potência média (ataque) & 3,48~W & \textbf{5,06~W} \\
\quad Energia total (banda)   & 4.233--6.454~J & 6.437--9.973~J \\
\bottomrule
\end{tabular}
\end{table}

O resultado central da comparação é que acrescentar a Fase~2 ao nó de borda tem
um custo real, e não desprezível, sob carga volumétrica. Enquanto o IDS binário
mantém a CPU em 11,8\% na média (pico de 44\%), o pipeline hierárquico a leva a
40,1\% na média (pico de 74\%), porque a segunda inferência roda em série sobre
cada alerta. Esse custo de processamento se traduz em latência e energia. A
latência mediana de classificação sobe de 120~s no binário para 352~s no
multiclasse, e a potência média durante o ataque de 3,48~W para 5,06~W, cerca de
45\% a mais. A memória cresce pouco entre os dois pipelines (RAM máxima de 2.528
para 2.663~MB), mas o nível absoluto, em torno de 2,6~GB, é ele próprio elevado:
vem da fila que armazena o \textit{backlog} de fluxos ainda não classificados,
com pico de quase meio milhão de fluxos. Rotular o tipo de ataque para triagem é,
portanto, viável no dispositivo de borda em regime normal, mas sob \textit{flood}
o custo da Fase~2 desloca de forma sensível a latência e a energia.

O volume de fluxos classificados difere pouco entre as sessões (831.154 no binário
contra 786.821 no multiclasse). As duas geram \textit{flood} independente, cuja
intensidade varia alguns por cento entre execuções, e o multiclasse, mais lento,
encerrou com um \textit{backlog} residual de 26.928 fluxos que não chegaram a ser
classificados dentro do tempo de encerramento. Descontado esse resíduo, as duas
sessões observam volumes equivalentes, a aproximadamente 1.050 fluxos de ataque
por segundo. A escolha entre os dois pipelines é orientada pelo uso: onde se
precisa de alta frequência de alertas brutos, por exemplo para alimentar um
sistema de gestão de eventos e informações de segurança (\textit{Security
Information and Event Management}, SIEM), o binário basta, e é mais barato. Onde se
precisa de rótulo de tipo para triagem, o multiclasse o agrega ao custo de CPU,
latência e energia acima.

A comparação confirma, lado a lado, os dois comportamentos de detecção já
discutidos. A precisão por fluxo é alta em ambos (98,4\% e 97,0\%, janela $+2$~s),
e a revocação binária, em torno de 64\%, concentra as perdas nos \textit{floods}
UDP de pacote único e nos ataques de baixo volume, enquanto as inundações
volumétricas de sessão TCP (\textit{spoofing}, DoS) são detectadas de forma
robusta. O FPR aparentemente alto não mede especificidade aqui, mas o denominador
diminuto de tráfego benigno sob captura saturada de \textit{flood}
(Seção~\ref{subsec:res_bin}).

Um limite de capacidade condiciona todos esses números
(Seção~\ref{subsec:res_artefatos}): a ~1.050 fluxos/s o nó ARM não classifica em
tempo real, e o IDS enfileira o tráfego para não descartá-lo, com pico de 456.681
a 495.051 fluxos na fila e latência mediana de 180 a 437~s. A detecção é avaliada
por fluxo justamente porque esse atraso desloca \emph{quando} o alerta aparece, não
\emph{qual} fluxo é. É um teto de vazão do nó de borda, retomado como limitação e
trabalho futuro na conclusão.

O consumo de energia da Tabela~\ref{tab:comparison} é uma estimativa, não uma
medição direta, porque o VIM~4 não expõe interface RAPL, INA226 ou
\texttt{hwmon}. A potência é modelada linearmente pela utilização de CPU,
\begin{equation}
P(t) = P_{idle} + (P_{max} - P_{idle}) \times \text{CPU\%}(t),
\label{eq:power}
\end{equation}
com \textit{endpoints} calibrados: \texttt{calibrate\_power.py} mede a utilização
de CPU em repouso (0,1\%) e sob carga total (\texttt{stress}, 99,9\%),
registrando como proveniência a temperatura (de 49 a 64~°C) e a frequência por
\textit{cluster} (A73 de 1800 a 2208~MHz; A53 de 500 a 2016~MHz), e ancora
$P_{idle} = 2{,}5$~W e $P_{max} = 9{,}0$~W no envelope de potência
\textit{board-level} do Khadas VIM~4 / A311D2 (USB-C PD). Os scripts do IDS
registram \texttt{power\_w}, temperatura e o escalonamento dinâmico de tensão e
frequência (\textit{Dynamic Voltage and Frequency Scaling}, DVFS) a cada amostra
($\sim$3~s), e a potência é integrada sobre essa série, separando ataque e
repouso pelas janelas de \textit{ground truth}. A incerteza é reportada como uma
banda de sensibilidade ($P_{idle} \in [2{,}0;3{,}0]$~W, $P_{max} \in
[7{,}0;11{,}0]$~W): de 4.233 a 6.454~J no binário e de 6.437 a 9.973~J no
multiclasse. O que importa é a comparação relativa: o consumo durante ataques
sobe de $\approx$3,5~W no binário para $\approx$5,1~W no multiclasse, coerente com
o custo de CPU da segunda inferência. Uma medição absoluta exigiria um medidor de
hardware externo, apontado como trabalho futuro.

% ============================================================================
\section{CONCLUSÕES}
\label{sec:conclusoes}

Este projeto apresentou um IDS hierárquico de duas fases para ambientes IoT em
redes 6G com MEC, cobrindo desde a construção do pipeline de dados até a
implantação e a validação no \textit{testbed} sobre a placa VIM~4, o que cumpre
o objetivo geral e os seis primeiros objetivos específicos da
Seção~\ref{sec:objetivos}; o sétimo, a terceira fase de agrupamento, foi
explorado apenas \textit{offline} (Seção~\ref{subsubsec:fase3}). A contribuição
central não é apenas o modelo, mas a caracterização empírica do que um IDS
hierárquico exige e entrega num nó de borda ARM real sob ataque volumétrico,
incluindo o ponto em que as \textit{features} por fluxo deixam de ser suficientes
e o teto de vazão do dispositivo. Em validação com o limiar de implantação fixado
em 0,9, sob métrica \emph{por fluxo} (cada fluxo rotulado pela janela em que
começou, imune ao atraso de classificação), o IDS binário alcançou precisão de
$\approx$98\% e revocação de $\approx$64\% (F1 de $\approx$78\%), mais forte nas
inundações volumétricas de sessão TCP, com \textit{spoofing} a 94\% e DoS a 92\%.
A revocação reflete o limiar de 0,9: os \textit{floods} UDP de pacote único, sem
estado bidirecional, ficam logo abaixo do limiar. Sobre os fluxos sinalizados, a
Fase~2 classificou corretamente o reconhecimento de rede (F1 = 73,7\%), mas
converge as demais classes para \texttt{dos}.

Acrescentar a Fase~2 ao nó de borda tem um custo real sob carga: leva a CPU média
de 12\% a 40\% e a potência de ataque de 3,5 para 5,1~W, além de triplicar a
latência de classificação (de 120 para 352~s de mediana). Rotular o tipo de ataque
para triagem é viável em regime normal, mas sob \textit{flood} esse custo não é
desprezível. A validação tornou explícito o teto da abordagem atual em duas
frentes. Na detecção de tipo, a Fase~2 converge para \texttt{dos} as classes cujo
comportamento ao vivo é, de fato, uma inundação de fluxos curtos e de alta taxa,
sobretudo o \texttt{spoofing}, que domina 119.873 fluxos. Para o \texttt{spoofing},
esse colapso é inclusive a resposta correta. Essa fronteira é um resultado, não
uma falha do pipeline: separar tais ataques exige \textit{features} de estado
agregado entre fluxos, como contadores de falha de autenticação e diversidade de
IPs de origem, que as 55 \textit{features} por fluxo não capturam. Na vazão, a
~1.050 fluxos/s o VIM~4 não classifica os fluxos em tempo real: o IDS enfileira o
tráfego para não descartá-lo, com pico de quase meio milhão de fluxos e latência
mediana de 180~s (binário) a 437~s (multiclasse). Soma-se a isso o desvio entre o
tráfego de treino e o de produção, que induz falsos positivos mitigados pela
inclusão de tráfego benigno real no treino e pela guarda de pacote mínimo. Toda a
validação foi
consolidada em scripts reprodutíveis, executáveis com um único comando, e o
projeto completo, incluindo os modelos treinados, está disponível em repositório
aberto\footnote{\url{https://github.com/luiz-linkezio/A-Hierarchical-Solution-based-on-MEC-and-IoT-for-Cyberattack-Detection-in-6G-Networks}};
a ferramenta \textit{netflower} é mantida em repositório próprio
(Seção~\ref{subsec:pipeline}).

Esses limites orientam os trabalhos futuros. O mais consequente é elevar o teto de
vazão do nó de borda, para que a classificação acompanhe o \textit{flood} em tempo
real sem enfileirar minutos de tráfego, seja por amostragem de fluxos sob carga,
por paralelismo de inferência em mais de um núcleo, ou por captura e extração mais
eficientes. Em seguida, cabe enriquecer o espaço de \textit{features} com estado
agregado entre fluxos, capaz de separar força bruta, \textit{spoofing} e MITM da
classe \texttt{dos}, acompanhado de um ambiente de avaliação com serviços e
vítimas reais que reduza o desvio entre treino e produção. Cabe ainda reduzir o
\texttt{idle\_timeout} para capturar ataques de curta duração como o \textit{scan}
UDP do \textit{recon}, avaliar um limiar de decisão mais baixo para recuperar os
\textit{floods} UDP de pacote único que o limiar de 0,9 descarta, substituir a
estimativa de energia por medição direta com um medidor de hardware externo e
integrar ao IDS a terceira fase de agrupamento explorada \textit{offline}, para
sinalizar tráfego que não se encaixa nas classes conhecidas. Por fim, como as
métricas reportadas derivam de uma única execução de cada sessão de avaliação,
cabe submetê-las a análise
estatística: repetir o protocolo experimental múltiplas vezes para estimar a
variabilidade das métricas de detecção e de custo, reportá-las com intervalos
de confiança e verificar, por testes de significância, se as diferenças
observadas entre os pipelines binário e multiclasse se sustentam entre
execuções.

% ============================================================================
% 8. REFERÊNCIAS BIBLIOGRÁFICAS (seção numerada via tocbibind, ver preâmbulo)
% ============================================================================
\bibliographystyle{sbc}
\bibliography{referencias}

% ============================================================================
\section{PRODUÇÃO BIBLIOGRÁFICA ADVINDA DA EXECUÇÃO DO PROJETO}
\label{sec:producao}

\begin{itemize}
  \item \textbf{Artigo científico completo} (em português, formato da Sociedade
    Brasileira de Computação, SBC): \textsc{Silva}, L. H. B. A.; \textsc{Souza},
    C. B. B.; \textsc{Balieiro}, A. M. \textit{Uma Solução Hierárquica para
    Detecção de Ciberataques em Ambientes IoT baseados em Redes MEC 6G}, 2026.
    O artigo consolida os resultados deste plano de trabalho e constituiu o
    Trabalho de Conclusão de Curso (TCC) do estudante no curso de Sistemas de
    Informação do Centro de Informática (CIn) da UFPE; até a data deste
    relatório, não havia sido aceito em evento ou periódico.
  \item \textbf{Versão resumida em inglês}, \textit{A Hierarchical Attack
    Detection Solution for 6G MEC-based IoT Networks}, preparada no formato
    IEEE para submissão a conferência internacional, também sem aceite até a
    data deste relatório.
  \item \textbf{Software} \textit{netflower}: extrator de fluxos de rede para
    dispositivos de borda, publicado no Python Package Index
    (\url{https://pypi.org/project/netflower/}) sob licença MIT, com
    código-fonte em \url{https://github.com/luiz-linkezio/netflower}.
  \item \textbf{Repositório aberto do projeto}, com os notebooks do pipeline de
    dados e de treinamento, os scripts do IDS e da validação automatizada, os
    modelos treinados e os relatórios de experimentos:
    \url{https://github.com/luiz-linkezio/A-Hierarchical-Solution-based-on-MEC-and-IoT-for-Cyberattack-Detection-in-6G-Networks}.
\end{itemize}

% ============================================================================
\section{ATIVIDADES PARALELAS DESENVOLVIDAS PELO ESTUDANTE}
\label{sec:paralelas}

Em paralelo ao plano de trabalho, o estudante elaborou seu Trabalho de
Conclusão de Curso no curso de Sistemas de Informação do CIn/UFPE, cujo tema
coincide com o desta pesquisa e cujo texto corresponde ao artigo listado na
Seção~\ref{sec:producao}. Também desenvolveu e publicou como software livre a
ferramenta \textit{netflower}, mantida em repositório próprio, cuja utilidade
extrapola este projeto por converter capturas de rede em fluxos com o conjunto
de \textit{features} do CICFlowMeter em qualquer máquina Linux, inclusive
dispositivos ARM.

% ============================================================================
\section{DIFICULDADES ENCONTRADAS NO DESENVOLVIMENTO DO PROJETO}
\label{sec:dificuldades}

As principais dificuldades surgiram na transição do ambiente \textit{offline}
para a operação real no nó de borda e estão documentadas, junto com as
soluções adotadas, no repositório do projeto.

A primeira foi a escala dos dados. As capturas brutas dos datasets CIC somam
cerca de 1~TB de PCAPs, inviáveis de manter em memória ou de processar com
ferramentas convencionais de extração de fluxos em tempo razoável, sobretudo
com vistas à execução posterior em um dispositivo ARM. A solução foi
desenvolver o \textit{netflower}, com processamento paralelo e memória
constante por fluxo, e consolidar os fluxos em um banco SQLite de cerca de
190~GB, inserido em blocos e indexado por rótulo, o que permitiu amostragem
balanceada em hardware convencional.

A segunda foi o desvio de distribuição entre o tráfego de treino e o de
produção. Modelos treinados apenas com o tráfego benigno \textit{scripted} dos
datasets classificavam como ataque o tráfego interativo real do laboratório
(sessões SSH e HTTP de curta duração), e um artefato do extrator (micro-fluxos
de encerramento TCP) gerava falsos positivos sistemáticos. Foi necessário
coletar dezenas de horas de tráfego benigno real para o treinamento, elevar o
limiar de implantação para 0,9 e adicionar uma guarda de pacote mínimo no IDS.
Parte do desvio persiste e explica, por exemplo, o \texttt{web} ao vivo ser
classificado como benigno, o que só uma avaliação com serviços e vítimas reais
poderia dissipar.

A terceira foi a própria execução dos experimentos no nó ARM. Sob
\textit{floods} de alta taxa, o nó não classifica os fluxos em tempo real, o que
desloca os alertas em minutos e motivou tanto a métrica \emph{por fluxo} (rotulando
cada fluxo pelo seu início, e não pela emissão) quanto a separação entre captura e
inferência, para que a inferência não bloqueie a captura e o \textit{kernel} não
descarte pacotes. Além disso, o IDS não encerrava com SIGINT, porque a captura ao
vivo sequestra o sinal, e contaminava a sessão seguinte; a captura PCAP no
atacante gerava arquivos de vários gigabytes sem valor analítico; e o ataque de
\textit{spoofing} inicialmente não emitia pacotes. Cada problema exigiu uma
correção no protocolo ou nos scripts: intervalos ociosos de 120~s entre ataques,
métrica por fluxo rotulada por \texttt{flow\_ts} (com janela estrita e calibração
$+2$~s), reinstalação dos tratadores de sinal com escalada INT/TERM/KILL e
carência de encerramento suficiente para o IDS drenar a fila, e desativação da
captura PCAP. O gerador também exigiu ajustes para que todos os ataques fossem
avaliáveis: a força bruta passou a usar vários usuários e listas maiores, e o MITM
a encaminhar tráfego TCP e UDP pela rota envenenada, já que o \textit{netflower}
converte em fluxo apenas TCP e UDP, e não o ICMP do \textit{ping} usado antes.

Por fim, a VIM~4 não expõe sensor de potência por software (RAPL, INA226 ou
\texttt{hwmon}), de modo que o consumo teve de ser estimado por um modelo linear
na utilização de CPU, calibrado no dispositivo e reportado com banda de incerteza;
uma medição absoluta exigiria um medidor de hardware externo. E a terceira fase da
arquitetura (agrupamento de fluxos de baixa confiança) ficou restrita a uma
exploração \textit{offline}, não integrada ao IDS implantado.

\vspace{0.7cm}
\noindent
\begin{tabular}{@{}p{7.4cm}@{\hspace{1.2cm}}p{7.4cm}@{}}
\rule{7.4cm}{0.4pt} & \rule{7.4cm}{0.4pt} \\
Data e assinatura do(a) orientador(a) & Data e assinatura do(a) estudante \\
\end{tabular}

\end{document}
